\documentclass[11pt,oneside]{book}
\usepackage[margin=1in]{geometry}
\usepackage{amsmath,amssymb,amsthm}
\usepackage{mathrsfs}
\usepackage{enumitem}
\usepackage{hyperref}
\usepackage{xcolor}
\usepackage{listings}
\usepackage{tikz}
\usepackage{epigraph}
\usepackage{tocloft}

\theoremstyle{plain}
\newtheorem{theorem}{Theorem}[chapter]
\newtheorem{lemma}[theorem]{Lemma}
\newtheorem{proposition}[theorem]{Proposition}
\newtheorem{corollary}[theorem]{Corollary}

\theoremstyle{definition}
\newtheorem{definition}[theorem]{Definition}
\newtheorem{example}[theorem]{Example}
\newtheorem{axiom_def}[theorem]{Axiom}

\theoremstyle{remark}
\newtheorem{remark}[theorem]{Remark}
\newtheorem{interpretation}[theorem]{Interpretation}

\lstdefinelanguage{Lean4}{
  morekeywords={theorem,lemma,def,noncomputable,class,instance,where,
    variable,open,namespace,section,end,import,structure,inductive,
    match,with,fun,let,in,if,then,else,do,return,have,show,by,
    exact,rfl,simp,linarith,nlinarith,ring,intro,apply,rw,unfold,
    calc,sorry,Prop,Type,Sort,axiom,extends,deriving},
  sensitive=true,
  morecomment=[l]{--},
  morecomment=[s]{/-}{-/},
  morestring=[b]",
}

\lstset{
  language=Lean4,
  basicstyle=\ttfamily\small,
  keywordstyle=\color{blue}\bfseries,
  commentstyle=\color{gray}\itshape,
  stringstyle=\color{red},
  breaklines=true,
  frame=single,
  backgroundcolor=\color{gray!5},
  xleftmargin=1em,
  xrightmargin=1em,
  literate={→}{{$\to$}}1 {←}{{$\leftarrow$}}1 {↦}{{$\mapsto$}}1
           {∀}{{$\forall$}}1 {∃}{{$\exists$}}1 {λ}{{$\lambda$}}1
           {≤}{{$\leq$}}1 {≥}{{$\geq$}}1 {≠}{{$\neq$}}1
           {∧}{{$\wedge$}}1 {∨}{{$\vee$}}1 {¬}{{$\neg$}}1
           {∈}{{$\in$}}1 {∉}{{$\notin$}}1 {⊆}{{$\subseteq$}}1
           {ℕ}{{$\mathbb{N}$}}1 {ℤ}{{$\mathbb{Z}$}}1 {ℝ}{{$\mathbb{R}$}}1
           {∘}{{$\circ$}}1 {×}{{$\times$}}1 {⟨}{{$\langle$}}1 {⟩}{{$\rangle$}}1
           {↔}{{$\leftrightarrow$}}1 {·}{{$\cdot$}}1 {ε}{{$\varepsilon$}}1
           {∩}{{$\cap$}}1 {₀}{{$_0$}}1,
}

\newcommand{\rs}{\mathrm{rs}}
\newcommand{\emerge}{\kappa}
\newcommand{\compose}{\circ}
\newcommand{\void}{\varepsilon}
\newcommand{\IS}{\mathcal{I}}
\newcommand{\R}{\mathbb{R}}
\newcommand{\N}{\mathbb{N}}

\title{\Huge\bfseries The Math of Ideas\\[0.3em]
\LARGE Volume IV-B: Paradoxes of Signs and Culture\\[0.5em]
\Large Counter-Intuitive Theorems in Semiotics, Narrative, and Aesthetics}
\author{A Formal Treatment with Machine-Verified Proofs\\[0.3em]
\normalsize\textit{Sequel to Volume~IV: Signs and Culture}}
\date{}

\begin{document}
\maketitle

% ================================================================
% PREFACE
% ================================================================
\frontmatter

\chapter*{Preface}
\addcontentsline{toc}{chapter}{Preface}

\epigraph{Everything we hear is an opinion, not a fact.
Everything we see is a perspective, not the truth.}{Marcus Aurelius, \textit{Meditations}}

This volume is the sequel to \textit{Volume~IV: Signs and Culture} and
the most philosophically provocative in the series.  Where Volume~IV
established the \emph{positive} theory---how signs work, how narratives
compose, how translation maps meaning---this volume develops the
\emph{negative} theory: the \textbf{paradoxes, impossibilities, and
counter-intuitive constraints} that the axioms of \textsc{IdeaticSpace}
impose on signs and culture.

The results here are \emph{surprising}.  They are not what a naive
reader would expect.  Some are outright paradoxical:

\begin{itemize}[leftmargin=2em]
\item \textbf{Adding beauty can destroy beauty} (the Kitsch Theorem,
  Chapter~4): the total aesthetic impact of a combination can be
  \emph{less} than either part alone, because emergence can be negative.

\item \textbf{Every metaphor necessarily distorts} (Chapter~6): there
  is no such thing as a ``perfect'' metaphor.  Every non-void vehicle
  introduces distortion for some observer.

\item \textbf{Translation loss is inevitable} (Chapter~2): any
  non-identity translation that changes self-resonance must introduce
  distortion.  You cannot translate ``approximately'' without
  ``approximately'' losing meaning.

\item \textbf{Novelty gets old} (Chapter~3): defamiliarization has
  diminishing returns.  The $n$-th weird thing you do is constrained by
  the cocycle condition to relate algebraically to the $(n{-}1)$-th.

\item \textbf{Order dominates content} (Chapter~5): reordering the
  elements of a narrative can change its meaning more than replacing
  an element entirely, because order sensitivity is \emph{purely}
  an emergence effect.

\item \textbf{Signs must drift} (Chapter~7): every non-linear sign
  necessarily changes meaning under recontextualization.  Stability
  requires triviality.

\item \textbf{Imitation creates originality} (Chapter~13): parody
  is always at least as ``substantial'' as its source, and the
  emergence at the point of parody is genuinely new meaning.
\end{itemize}

Every theorem in this volume is \textbf{machine-verified} in Lean~4,
building on the \texttt{IdeaticSpace3} axioms from
\texttt{IDT\_v3\_Foundations.lean}.  The proofs use \textbf{zero
sorries}---every step is checked by the Lean kernel.  The relevant
Lean file is \texttt{IDT\_v3\_DeepSemiotics.lean}.

\vspace{1em}

\noindent\textbf{Note on axioms.}  All results follow from the eight
axioms of \texttt{IdeaticSpace3}: associative composition with identity
(A1--A3), resonance annihilation by void (R1--R2), non-negative
self-resonance (E1), nondegeneracy (E2), bounded emergence (E3), and
compositional enrichment (E4).  No additional axioms are introduced.

\tableofcontents

% ================================================================
\mainmatter

% ══════════════════════════════════════════════════════════════════════
% Volume IV-B, Chapters 1--5
% ══════════════════════════════════════════════════════════════════════

% ══════════════════════════════════════════════════════════════════════
% Chapter 1: The Decay of Meaning
% ══════════════════════════════════════════════════════════════════════

\chapter{The Decay of Meaning: Why Reinterpretation Always Loses}
\label{ch:decay}

\epigraph{``A word is dead / When it is said, / Some say. /
I say it just / Begins to live / That day.''}{Emily Dickinson}

% ── Historical Preamble ──────────────────────────────────────────────

The childhood game of ``telephone''---in which a message is whispered
from ear to ear around a circle, arriving at the end comically
garbled---is not merely a game.  It is a \emph{theorem}.

In this chapter, we prove that every act of reinterpretation---every
composition of a sign with a new context---introduces an irreversible
distortion.  The distortion is not random noise; it has a precise
algebraic structure.  It decomposes into a \emph{bias term} (the
context's own resonance with the observer) and an \emph{emergence term}
(the nonlinear interaction between sign and context).  Only when both
are zero---only when the context is silence---is the reinterpretation
lossless.

This result overturns the naive hope that ``careful'' reinterpretation
can preserve meaning perfectly.  It cannot.  Every reading leaves a mark
on the text, and the marks accumulate.

\section{Sign Chains and Reinterpretation}

We begin by formalizing the notion of a \emph{sign chain}: a sequence
of reinterpretations, each modeled as composition with a new context.

\begin{definition}[Sign Chain]
\label{def:sign-chain}
Let $a \in \IS$ be a sign and $c_1, c_2, \ldots, c_n \in \IS$ be
a sequence of contexts.  The \emph{$n$-step sign chain} is defined
recursively:
\[
  \mathrm{chain}(a, []) = a, \qquad
  \mathrm{chain}(a, c :: \mathrm{rest}) = \mathrm{chain}(a \compose c, \mathrm{rest}).
\]
\end{definition}

\begin{lstlisting}
def signChain (a : I) : List I → I
  | [] => a
  | c :: rest => signChain (compose a c) rest
\end{lstlisting}

The definition is straightforward, but its consequences are profound.

\begin{theorem}[Sign Weight Growth]
\label{thm:sign-weight-growth}
For any sign $a$ and context $c$,
\[
  \rs(a \compose c, a \compose c) \geq \rs(a, a).
\]
Every reinterpretation increases or preserves the sign's weight.
\end{theorem}

\begin{proof}
This is an immediate consequence of axiom E4 (compositional enrichment):
$\rs(a \compose b, a \compose b) \geq \rs(a, a)$ for all $a, b$.
\end{proof}

\begin{remark}[Why this is surprising]
At first glance, this seems like good news: reinterpretation makes signs
``heavier.''  But weight measures \emph{quantity of presence}, not
\emph{fidelity to the original meaning}.  A sign can grow heavier while
drifting arbitrarily far from its original resonance profile.  The sign
gets more ``substantial'' while becoming less ``faithful.''  This is the
paradox of tradition: cultural transmission accumulates weight while
potentially distorting meaning.
\end{remark}

\section{The Reinterpretation Gap}

\begin{definition}[Reinterpretation Gap]
The \emph{reinterpretation gap} measures how much one step of
reinterpretation changes the resonance of a sign with an observer:
\[
  \mathrm{gap}(a, c, \mathrm{obs}) := \rs(a \compose c, \mathrm{obs}) - \rs(a, \mathrm{obs}).
\]
\end{definition}

\begin{lstlisting}
noncomputable def reinterpretationGap (a c observer : I) : R :=
  rs (compose a c) observer - rs a observer
\end{lstlisting}

\begin{theorem}[Gap Decomposition]
\label{thm:gap-decomp}
The reinterpretation gap decomposes as:
\[
  \mathrm{gap}(a, c, \mathrm{obs}) = \rs(c, \mathrm{obs}) + \emerge(a, c, \mathrm{obs}).
\]
\end{theorem}

\begin{proof}
By the resonance composition equation $\rs(a \compose c, \mathrm{obs})
= \rs(a, \mathrm{obs}) + \rs(c, \mathrm{obs}) + \emerge(a, c, \mathrm{obs})$,
subtraction gives the result immediately.
\end{proof}

\begin{lstlisting}
theorem reinterpretationGap_decomp (a c observer : I) :
    reinterpretationGap a c observer =
    rs c observer + emergence a c observer := by
  unfold reinterpretationGap
  have h := rs_compose_eq a c observer
  linarith
\end{lstlisting}

\begin{interpretation}[The two channels of distortion]
The gap has two independent components:
\begin{enumerate}
\item The \textbf{bias term} $\rs(c, \mathrm{obs})$: the context's own
  resonance with the observer.  This is the ``literal'' contribution of
  the context---what it brings to the table independently of the sign.
\item The \textbf{emergence term} $\emerge(a, c, \mathrm{obs})$: the
  nonlinear interaction between sign and context.  This is the ``creative''
  distortion---the new meaning that arises from the composition.
\end{enumerate}
Both terms are generally nonzero.  To achieve zero gap, \emph{both}
must vanish simultaneously.  This is a much stronger condition than
either alone.
\end{interpretation}

\section{Two-Step Distortion: The Telephone Theorem}

\begin{theorem}[Two-Step Distortion]
\label{thm:two-step}
When a sign is reinterpreted through two successive contexts
$c_1, c_2$, the total gap equals the sum of individual gaps
\emph{plus} a cross-emergence term:
\[
  \rs(a \compose c_1 \compose c_2, \mathrm{obs}) - \rs(a, \mathrm{obs})
  = \mathrm{gap}(a, c_1, \mathrm{obs}) +
    \mathrm{gap}(a \compose c_1, c_2, \mathrm{obs}).
\]
\end{theorem}

\begin{proof}
This is a telescoping identity:
\begin{align*}
  &\rs((a \compose c_1) \compose c_2, \mathrm{obs}) - \rs(a, \mathrm{obs}) \\
  &= [\rs((a \compose c_1) \compose c_2, \mathrm{obs})
    - \rs(a \compose c_1, \mathrm{obs})] +
    [\rs(a \compose c_1, \mathrm{obs}) - \rs(a, \mathrm{obs})] \\
  &= \mathrm{gap}(a \compose c_1, c_2, \mathrm{obs}) +
    \mathrm{gap}(a, c_1, \mathrm{obs}).
\end{align*}
\end{proof}

\begin{lstlisting}
theorem two_step_distortion (a c1 c2 observer : I) :
    reinterpretationGap a c1 observer +
    reinterpretationGap (compose a c1) c2 observer =
    rs (compose (compose a c1) c2) observer - rs a observer := by
  unfold reinterpretationGap; ring
\end{lstlisting}

\begin{remark}[Why this is surprising]
The naive expectation is that two ``small'' reinterpretations should
produce a ``small'' total distortion.  But each gap decomposes into
bias + emergence, and the \emph{emergence terms are generally correlated}
via the cocycle condition.  Two individually small gaps can produce a
large total distortion if the cross-emergence aligns.  Distortion is
not additive---it is \emph{super-additive} when emergence cooperates.
\end{remark}

\section{The Only Lossless Reinterpretation}

\begin{theorem}[Void Preserves Perfectly]
\label{thm:void-preserves}
The only context that produces zero reinterpretation gap for all
observers is the void:
\[
  \mathrm{gap}(a, \void, \mathrm{obs}) = 0 \quad \text{for all } a, \mathrm{obs}.
\]
\end{theorem}

\begin{proof}
By axiom A3, $a \compose \void = a$, so
$\rs(a \compose \void, \mathrm{obs}) = \rs(a, \mathrm{obs})$
and the gap is zero.
\end{proof}

\begin{lstlisting}
theorem void_reinterpretation_zero (a observer : I) :
    reinterpretationGap a (void : I) observer = 0 := by
  unfold reinterpretationGap; simp [void_right']
\end{lstlisting}

\begin{interpretation}[Silence is the only perfect reading]
This theorem says that the \emph{only} way to ``reinterpret'' a text
without changing it is to add nothing.  Every substantive reading---every
context that is not pure silence---necessarily changes what the text
means to some observer.  This is the formal version of reader-response
theory's central claim: the reader always contributes.  And the only
reader who contributes nothing is the one who doesn't read.
\end{interpretation}

\section{Non-Void Signs Leave Permanent Marks}

\begin{theorem}[Irreversibility of Reading]
\label{thm:reading-irreversible}
If $a \neq \void$, then for any context $c$:
\[
  \rs(a \compose c, a \compose c) > 0.
\]
The reinterpreted sign always has strictly positive weight.
\end{theorem}

\begin{proof}
Since $a \neq \void$, we have $\rs(a, a) > 0$ by the nondegeneracy axiom.
By compositional enrichment, $\rs(a \compose c, a \compose c) \geq \rs(a, a) > 0$.
\end{proof}

\begin{lstlisting}
theorem nonvoid_reinterpretation_positive (a c : I)
    (ha : a ≠ void) :
    rs (compose a c) (compose a c) > 0 := by
  have h1 := compose_enriches' a c
  have h2 := rs_self_pos a ha
  linarith
\end{lstlisting}

\begin{remark}[Why this is surprising]
This means that once an idea enters the cultural conversation, it can
never be fully erased.  Composition with any context---even a hostile
one, even an attempt at refutation---produces something with positive
weight.  Censorship, in the ideatic-space sense, is impossible: engaging
with an idea to suppress it only adds to its weight.  The only way to
truly silence an idea is to never engage with it at all (compose with
void).
\end{remark}


% ══════════════════════════════════════════════════════════════════════
% Chapter 2: The Impossibility of Perfect Translation
% ══════════════════════════════════════════════════════════════════════

\chapter{The Impossibility of Perfect Translation}
\label{ch:translation}

\epigraph{``Traduttore, traditore.''\\(Translator, traitor.)}{Italian proverb}

Translation has been called ``the art of failure'' (Umberto Eco).  In this
chapter, we prove that this is not merely a literary sentiment but a
\emph{mathematical theorem}.  Any translation that is not the identity
must introduce distortion.  The distortion is measurable, it accumulates
through chains of retranslation, and it has a precise lower bound.

\section{Translation Loss: The Basic Inequality}

\begin{definition}[Translation Loss]
For a translation map $\varphi : \IS \to \IS$, the \emph{translation
loss} for a pair $(a, b)$ is:
\[
  \mathrm{loss}(\varphi, a, b) := |\rs(\varphi(a), \varphi(b)) - \rs(a, b)|.
\]
\end{definition}

\begin{lstlisting}
noncomputable def translationLoss (phi : I → I) (a b : I) : R :=
  |rs (phi a) (phi b) - rs a b|
\end{lstlisting}

\begin{theorem}[Faithful Translations Have Zero Loss]
\label{thm:faithful-zero}
If $\varphi$ is faithful (preserves all resonances), then
$\mathrm{loss}(\varphi, a, b) = 0$ for all $a, b$.
\end{theorem}

\begin{proof}
If $\rs(\varphi(a), \varphi(b)) = \rs(a, b)$ for all $a, b$,
then the absolute value of their difference is zero.
\end{proof}

\begin{theorem}[The Self-Distortion Lower Bound]
\label{thm:self-distortion}
If a translation changes the self-resonance of even \emph{one} idea,
then the translation loss is strictly positive for that idea:
\[
  \rs(\varphi(a), \varphi(a)) \neq \rs(a, a)
  \implies
  \mathrm{loss}(\varphi, a, a) > 0.
\]
\end{theorem}

\begin{proof}
If $\rs(\varphi(a), \varphi(a)) \neq \rs(a, a)$, then their difference
is nonzero, so its absolute value is positive.
\end{proof}

\begin{lstlisting}
theorem translation_self_loss (phi : I → I) (a : I)
    (h : rs (phi a) (phi a) ≠ rs a a) :
    translationLoss phi a a > 0 := by
  unfold translationLoss
  exact abs_pos.mpr (sub_ne_zero.mpr h)
\end{lstlisting}

\begin{remark}[Why this is surprising]
The theorem says that you cannot translate ``approximately faithfully.''
If a translation preserves \emph{all} resonances except one self-resonance,
it is still unfaithful, and the unfaithfulness is \emph{detectable}:
the translation loss is strictly positive.  There is no ``good enough''
approximation---faithfulness is all or nothing.
\end{remark}

\section{Chain Translation: Distortion Accumulates}

\begin{theorem}[Chain Distortion Decomposition]
\label{thm:chain-decomp}
Composing two translations $\psi \circ \varphi$ produces distortion
that decomposes into two layers:
\[
  \rs(\psi(\varphi(a)), \psi(\varphi(b))) - \rs(a, b)
  = \underbrace{[\rs(\psi(\varphi(a)), \psi(\varphi(b)))
    - \rs(\varphi(a), \varphi(b))]}_{\text{second layer}} +
    \underbrace{[\rs(\varphi(a), \varphi(b)) - \rs(a, b)]}_{\text{first layer}}.
\]
\end{theorem}

\begin{proof}
This is a telescoping identity; it follows from algebraic rearrangement.
\end{proof}

\begin{lstlisting}
theorem chain_translation_loss_decomp (phi psi : I → I)
    (a b : I) :
    rs (psi (phi a)) (psi (phi b)) - rs a b =
    (rs (psi (phi a)) (psi (phi b)) - rs (phi a) (phi b)) +
    (rs (phi a) (phi b) - rs a b) := by ring
\end{lstlisting}

\begin{interpretation}[Why retranslation degrades]
Each layer of translation adds its own distortion.  Even if each individual
translation is ``small'' (low distortion), the total can be large because
the layers are \emph{independent}.  This is the formal reason that
translating a poem from Chinese to French to English produces a worse
result than translating directly from Chinese to English: each intermediate
step adds its own layer of loss, and the losses add.
\end{interpretation}

\section{Round-Trip Loss: The Myth of Perfect Back-Translation}

\begin{theorem}[Round-Trip Loss]
\label{thm:round-trip}
For any two translations $\varphi, \psi$, the round-trip distortion
satisfies:
\[
  \rs(\psi(\varphi(a)), \psi(\varphi(a))) - \rs(a, a)
  = \underbrace{[\rs(\psi(\varphi(a)), \psi(\varphi(a)))
    - \rs(\varphi(a), \varphi(a))]}_{\text{back-translation loss}} +
    \underbrace{[\rs(\varphi(a), \varphi(a)) - \rs(a, a)]}_{\text{forward loss}}.
\]
\end{theorem}

\begin{proof}
Again a telescoping identity.
\end{proof}

\begin{lstlisting}
theorem roundtrip_loss (phi psi : I → I) (a : I) :
    rs (psi (phi a)) (psi (phi a)) - rs a a =
    (rs (psi (phi a)) (psi (phi a)) - rs (phi a) (phi a)) +
    (rs (phi a) (phi a) - rs a a) := by ring
\end{lstlisting}

\begin{remark}[Why this is surprising]
The theorem destroys the comforting notion of ``back-translation as
quality check.''  Even if you translate and back-translate, the round-trip
distortion is the \emph{sum} of forward and backward losses.  Having a
``back-translation'' does not guarantee zero net loss.  The round trip
through Chinese $\to$ French $\to$ Chinese loses meaning in \emph{both}
directions.
\end{remark}

\section{Void-Preservation Is Not Sufficient}

\begin{theorem}[Void-Preserving Is Not Faithful]
\label{thm:void-not-faithful}
A void-preserving translation ($\varphi(\void) = \void$) can still
have arbitrarily high translation loss.  Void-preservation is necessary
but not sufficient for faithfulness.
\end{theorem}

\begin{proof}
We prove a disjunction: either $\varphi$ is faithful (zero loss everywhere),
or there exist $a, b$ with positive loss.  If $\varphi$ is not faithful,
there exist $a, b$ with $\rs(\varphi(a), \varphi(b)) \neq \rs(a, b)$,
giving positive loss.
\end{proof}

\begin{lstlisting}
theorem void_preserving_not_sufficient :
    ∀ phi : I → I, ds_voidPreserving phi →
    (ds_faithful phi ∨ ∃ a b, translationLoss phi a b > 0) := by
  intro phi _
  by_cases hf : ds_faithful phi
  · left; exact hf
  · right
    unfold ds_faithful at hf; push_neg at hf
    rcases hf with ⟨a, b, hab⟩
    exact ⟨a, b, by unfold translationLoss;
      exact abs_pos.mpr (sub_ne_zero.mpr hab)⟩
\end{lstlisting}


% ══════════════════════════════════════════════════════════════════════
% Chapter 3: Defamiliarization Fatigue
% ══════════════════════════════════════════════════════════════════════

\chapter{Defamiliarization Fatigue: Novelty Gets Old}
\label{ch:defamiliarization}

\epigraph{``Art exists that one may recover the sensation of life;
it exists to make one feel things, to make the stone stony.''}{Viktor Shklovsky, \textit{Art as Technique} (1917)}

Shklovsky's \emph{defamiliarization} (\textit{ostranenie})---the technique
of making the familiar strange---is central to modernist aesthetics.
But does defamiliarization have diminishing returns?  Can you keep
``making strange'' with undiminished effect?

In this chapter, we prove that the answer is \textbf{no}.  The cocycle
condition forces each successive defamiliarization to relate
algebraically to the previous one.  The artist cannot simply ``keep doing
the same weird thing'' and expect the same effect.  Novelty is
algebraically constrained.

\section{Iterated Defamiliarization}

\begin{definition}[Iterated Defamiliarization]
Applying the same defamiliarizing context $c$ to a sign $a$
repeatedly:
\[
  D^0(a, c) = a, \qquad D^{n+1}(a, c) = D^n(a, c) \compose c.
\]
\end{definition}

\begin{lstlisting}
def iteratedDefamiliarize (a ctx : I) : N → I
  | 0 => a
  | n + 1 => compose (iteratedDefamiliarize a ctx n) ctx
\end{lstlisting}

\begin{theorem}[Weight Grows Monotonically]
\label{thm:weight-mono}
Each iteration of defamiliarization strictly increases or preserves weight:
\[
  \rs(D^{n+1}(a,c), D^{n+1}(a,c)) \geq \rs(D^n(a,c), D^n(a,c)).
\]
\end{theorem}

\begin{proof}
$D^{n+1}(a,c) = D^n(a,c) \compose c$, so by axiom E4:
$\rs(D^n(a,c) \compose c, D^n(a,c) \compose c) \geq \rs(D^n(a,c), D^n(a,c))$.
\end{proof}

\begin{lstlisting}
theorem defamiliarize_weight_grows (a ctx : I) (n : N) :
    rs (iteratedDefamiliarize a ctx (n + 1))
       (iteratedDefamiliarize a ctx (n + 1)) ≥
    rs (iteratedDefamiliarize a ctx n)
       (iteratedDefamiliarize a ctx n) := by
  simp only [iteratedDefamiliarize]
  exact compose_enriches' (iteratedDefamiliarize a ctx n) ctx
\end{lstlisting}

\section{The Cocycle Constraint on Successive Novelty}

\begin{theorem}[Defamiliarization Cocycle]
\label{thm:defam-cocycle}
The emergence at the second defamiliarization is constrained by the
emergence at the first:
\[
  \emerge(a, c, \mathrm{obs}) + \emerge(a \compose c, c, \mathrm{obs})
  = \emerge(c, c, \mathrm{obs}) + \emerge(a, c \compose c, \mathrm{obs}).
\]
\end{theorem}

\begin{proof}
This is the cocycle condition with $b = c_1 = c$ (the same context
applied twice).
\end{proof}

\begin{lstlisting}
theorem defamiliarization_cocycle (a ctx obs : I) :
    emergence a ctx obs +
    emergence (compose a ctx) ctx obs =
    emergence ctx ctx obs +
    emergence a (compose ctx ctx) obs :=
  cocycle_condition a ctx ctx obs
\end{lstlisting}

\begin{remark}[Why this is surprising]
This is the key result.  It says that the emergence of the \emph{second}
application of context $c$ is \emph{not independent} of the first.
Rearranging:
\[
  \emerge(a \compose c, c, \mathrm{obs}) =
  \emerge(c, c, \mathrm{obs}) +
  \emerge(a, c \compose c, \mathrm{obs}) -
  \emerge(a, c, \mathrm{obs}).
\]
The second defamiliarization's effect is completely determined by three
quantities: (1) the self-emergence of $c$ (the context's own novelty),
(2) the emergence of $a$ with the doubled context $c \compose c$, and
(3) the emergence of the first application.  The artist cannot freely
choose the second novelty effect---it is constrained.

In particular, if $c$ has zero self-emergence ($\emerge(c,c,\mathrm{obs}) = 0$)
and $a$'s emergence with $c \compose c$ equals its emergence with $c$,
then the second application has \emph{exactly the same} effect as the
first.  And the third application is then constrained by the second.
The diminishing returns are structural, not accidental.
\end{remark}

\section{The Habituation Theorem}

\begin{theorem}[Habituation Constraint]
\label{thm:habituation}
After repeated defamiliarization with the same context, the emergence
at each step is fully determined by earlier steps:
\[
  \emerge(a \compose c, c, \mathrm{obs})
  = \emerge(c, c, \mathrm{obs}) + \emerge(a, c^2, \mathrm{obs})
    - \emerge(a, c, \mathrm{obs}).
\]
\end{theorem}

\begin{proof}
Rearranging the cocycle condition.
\end{proof}

\begin{lstlisting}
theorem habituation_constraint (a ctx obs : I) :
    emergence (compose a ctx) ctx obs =
    emergence ctx ctx obs +
    emergence a (compose ctx ctx) obs -
    emergence a ctx obs := by
  linarith [cocycle_condition a ctx ctx obs]
\end{lstlisting}

\begin{interpretation}[The algebra of artistic fatigue]
This theorem formalizes what every artist knows intuitively: you can't
keep doing the same trick.  The effect of the $n$-th weird thing is
algebraically determined by the $(n{-}1)$-th.  Each new iteration must
satisfy the cocycle constraint, which limits the freedom of the artist.

This is not merely a bound---it is an \emph{exact equation}.  The
emergence at step $n+1$ is a linear function of the emergence at step $n$,
plus corrections from the context's self-interaction and the
``accumulated'' emergence.  The artist's freedom is \emph{one degree of
freedom less} than naive independence would suggest.

In the language of cohomology, defamiliarization defines a 1-cochain,
and the cocycle condition says this cochain must be a cocycle.  The
artist is constrained to move within a cohomology class.
\end{interpretation}

\section{Void Defamiliarization Does Nothing}

\begin{theorem}[Void Defamiliarization Is Trivial]
\label{thm:void-defam}
Using silence as a defamiliarizing context has zero effect:
\[
  D^n(a, \void) = a \quad \text{for all } n.
\]
\end{theorem}

\begin{proof}
By induction on $n$.  Base case: $D^0(a, \void) = a$.
Inductive step: $D^{n+1}(a, \void) = D^n(a, \void) \compose \void
= a \compose \void = a$.
\end{proof}

\begin{lstlisting}
theorem void_defamiliarization (a obs : I) (n : N) :
    iteratedDefamiliarize a (void : I) n = a := by
  induction n with
  | zero => rfl
  | succ n ih => simp [iteratedDefamiliarize, ih, void_right']
\end{lstlisting}

\section{The Emergence Ratio Shrinks}

\begin{theorem}[Emergence-Weight Ratio Bound]
\label{thm:ratio-bound}
At each step, the emergence is bounded by
\[
  \emerge(a, c, \mathrm{obs})^2 \leq \rs(a \compose c, a \compose c) \cdot \rs(\mathrm{obs}, \mathrm{obs}).
\]
Since the left-hand side involves $\rs(a \compose c, a \compose c)$
which grows monotonically (Theorem~\ref{thm:weight-mono}), the
\emph{ratio} $\emerge^2 / \rs(D^n, D^n)$ can only decrease as the
weight grows.
\end{theorem}

\begin{proof}
Axiom E3 (emergence bounded).
\end{proof}

\begin{remark}[Relative novelty diminishes]
Even if the absolute emergence stays constant (a theoretical possibility),
the \emph{relative} novelty---emergence as a fraction of total weight---must
decrease.  This is the diminishing-returns theorem: each iteration of
defamiliarization contributes proportionally less to the total meaning
of the accumulated sign chain.  The $100$th repetition of ``make it
strange'' adds less \emph{relative} novelty than the first.
\end{remark}


% ══════════════════════════════════════════════════════════════════════
% Chapter 4: The Kitsch Theorem
% ══════════════════════════════════════════════════════════════════════

\chapter{The Kitsch Theorem: When Beauty Destroys Beauty}
\label{ch:kitsch}

\epigraph{``Kitsch causes two tears to flow in quick succession.
The first tear says: How nice to see children running on the grass!
The second tear says: How nice to be moved, together with all mankind,
by children running on the grass!''}{Milan Kundera, \textit{The Unbearable Lightness of Being}}

\section{Aesthetic Non-Monotonicity}

The most counter-intuitive result in our entire theory may be this:
\textbf{adding a beautiful element to a beautiful composition can make
the total aesthetic impact worse}.  In the language of economics, aesthetic
value is not a monotone function of its inputs.

\begin{definition}[Total Aesthetic Impact]
\[
  \mathrm{impact}(a, b, \mathrm{obs}) := \rs(a \compose b, \mathrm{obs}).
\]
\end{definition}

\begin{theorem}[The Kitsch Theorem]
\label{thm:kitsch}
\[
  \mathrm{impact}(a, b, \mathrm{obs})
  = \rs(a, \mathrm{obs}) + \rs(b, \mathrm{obs}) + \emerge(a, b, \mathrm{obs}).
\]
When $\emerge(a, b, \mathrm{obs}) < 0$, the total impact is \emph{less}
than the sum of individual impacts.
\end{theorem}

\begin{proof}
The resonance composition equation gives the decomposition directly.
\end{proof}

\begin{lstlisting}
theorem kitsch_theorem (a b observer : I) :
    totalAestheticImpact a b observer =
    rs a observer + rs b observer +
    emergence a b observer := by
  unfold totalAestheticImpact; rw [rs_compose_eq]
\end{lstlisting}

\begin{remark}[Why this is surprising]
This overturns the ``more is more'' intuition in aesthetics.  Consider:
\begin{itemize}
\item A beautifully written novel with an unnecessarily beautiful
  epilogue that undermines the ambiguity of the ending.
\item A perfectly composed photograph with an added Instagram filter
  that makes it look cheap.
\item A Baroque church with one too many gold cherubs.
\end{itemize}
In each case, the added element has positive resonance on its own, but
its \emph{emergence} with the existing work is negative.  The
composition clashes.  The whole is less than the sum of its parts.

This is the mathematical definition of \textbf{kitsch}: art where the
emergence is zero or negative, where addition is purely additive or
actively destructive.  The kitsch sensibility adds more beauty without
concern for how it interacts with what's already there.
\end{remark}

\section{The Over-Ornamentation Bound}

\begin{theorem}[Over-Ornamentation]
\label{thm:over-ornament}
If emergence is negative ($\emerge(a, b, \mathrm{obs}) < 0$), then
\[
  \mathrm{impact}(a, b, \mathrm{obs}) < \rs(a, \mathrm{obs}) + \rs(b, \mathrm{obs}).
\]
Adding $b$ makes the total impact \emph{strictly worse} than having
both parts separately.
\end{theorem}

\begin{proof}
Follows immediately from the Kitsch Theorem and the sign of the
emergence term.
\end{proof}

\begin{lstlisting}
theorem over_ornamentation (a b observer : I)
    (h : emergence a b observer < 0) :
    totalAestheticImpact a b observer <
    rs a observer + rs b observer := by
  unfold totalAestheticImpact
  linarith [rs_compose_eq a b observer]
\end{lstlisting}

\section{The Goldilocks Principle}

\begin{theorem}[The Goldilocks Principle]
\label{thm:goldilocks}
For any composition $a \compose b$:
\begin{enumerate}
\item If $\emerge(a, b, \mathrm{obs}) = 0$: impact equals the sum of parts (neutral).
\item If $\emerge(a, b, \mathrm{obs}) > 0$: impact exceeds the sum (synergy).
\item If $\emerge(a, b, \mathrm{obs}) < 0$: impact is less than the sum (clash).
\end{enumerate}
\end{theorem}

\begin{proof}
All three follow from the Kitsch Theorem by comparing signs.
\end{proof}

\begin{lstlisting}
theorem goldilocks_principle (a b obs : I) :
    (emergence a b obs = 0 →
      totalAestheticImpact a b obs = rs a obs + rs b obs) ∧
    (emergence a b obs > 0 →
      totalAestheticImpact a b obs > rs a obs + rs b obs) ∧
    (emergence a b obs < 0 →
      totalAestheticImpact a b obs < rs a obs + rs b obs) := by
  unfold totalAestheticImpact
  refine ⟨?_, ?_, ?_⟩ <;> intro h <;>
    linarith [rs_compose_eq a b obs]
\end{lstlisting}

\begin{interpretation}[The aesthetics of emergence]
The Goldilocks Principle says that aesthetic composition is not about
the quality of individual elements but about the \emph{emergence} between
them.  A mediocre element with positive emergence contributes more than
a brilliant element with negative emergence.  The art is in the
combination, not the parts.

This is the mathematical content of the old clich\'e ``less is more'':
removing an element with negative emergence \emph{increases} the total
impact.  The Goldilocks Principle gives the exact criterion for when
``more is more,'' ``more is less,'' and ``more is just more.''
\end{interpretation}

\section{Double Kitsch}

\begin{theorem}[Double Kitsch Is Predictable]
\label{thm:double-kitsch}
If $a$ is left-linear (kitsch), then
\[
  \mathrm{impact}(a, b, \mathrm{obs}) = \rs(a, \mathrm{obs}) + \rs(b, \mathrm{obs}).
\]
Kitsch produces exactly the sum of parts---no surprise, no synergy,
no clash.
\end{theorem}

\begin{proof}
Left-linearity means $\emerge(a, b, c) = 0$ for all $b, c$.
The Kitsch Theorem then gives the result.
\end{proof}

\begin{lstlisting}
theorem double_kitsch (a b observer : I)
    (ha : leftLinear a) :
    totalAestheticImpact a b observer =
    rs a observer + rs b observer := by
  unfold totalAestheticImpact
  rw [rs_compose_eq]; linarith [ha b observer]
\end{lstlisting}

\section{Weight Grows But Resonance Can Shrink}

\begin{theorem}[The Over-Production Paradox]
\label{thm:over-production}
The self-resonance (weight) of a composition always exceeds that of its
first element, but the resonance with a \emph{specific observer} can
decrease:
\[
  \rs(a \compose b, a \compose b) \geq \rs(a, a)
\]
but it is possible that $\rs(a \compose b, \mathrm{obs}) < \rs(a, \mathrm{obs})$.
\end{theorem}

\begin{proof}
The first inequality is axiom E4.  For the second, when
$\rs(b, \mathrm{obs}) + \emerge(a, b, \mathrm{obs}) < 0$,
we have $\rs(a \compose b, \mathrm{obs}) < \rs(a, \mathrm{obs})$.
\end{proof}

\begin{remark}[Why this is surprising]
This is the paradox of the over-produced film.  A film with a
\$200 million budget (high weight) can be less impactful to audiences
(less resonance with the observer) than a \$2 million indie film.
Weight measures \emph{quantity of stuff}; resonance measures
\emph{impact on a particular audience}.  They are different things,
and one can increase while the other decreases.
\end{remark}


% ══════════════════════════════════════════════════════════════════════
% Chapter 5: Narrative Order Dominates Content
% ══════════════════════════════════════════════════════════════════════

\chapter{Narrative Order Dominates Content}
\label{ch:order}

\epigraph{``The king died and then the queen died'' is a story.
``The king died, and then the queen died of grief'' is a plot.}{E.\ M.\ Forster, \textit{Aspects of the Novel}}

\section{The Order Sensitivity Theorem}

One of the most striking results in narrative theory: the \emph{order}
in which elements appear can change the meaning more than replacing
an element entirely.

\begin{definition}[Narrative Order Effect]
The \emph{order effect} of swapping $a$ and $b$ in a narrative:
\[
  \mathrm{order}(a, b, \mathrm{obs}) :=
  \rs(a \compose b, \mathrm{obs}) - \rs(b \compose a, \mathrm{obs}).
\]
\end{definition}

\begin{theorem}[Order Effect Equals Emergence Difference]
\label{thm:order-emergence}
\[
  \mathrm{order}(a, b, \mathrm{obs}) =
  \emerge(a, b, \mathrm{obs}) - \emerge(b, a, \mathrm{obs}).
\]
\end{theorem}

\begin{proof}
Expanding both resonances using the composition equation:
\begin{align*}
  \rs(a \compose b, \mathrm{obs}) &= \rs(a, \mathrm{obs}) + \rs(b, \mathrm{obs})
    + \emerge(a, b, \mathrm{obs}) \\
  \rs(b \compose a, \mathrm{obs}) &= \rs(b, \mathrm{obs}) + \rs(a, \mathrm{obs})
    + \emerge(b, a, \mathrm{obs}).
\end{align*}
Subtracting:
$\mathrm{order}(a, b, \mathrm{obs}) = \emerge(a, b, \mathrm{obs}) - \emerge(b, a, \mathrm{obs})$.
\end{proof}

\begin{lstlisting}
theorem narrativeOrderEffect_eq (a b observer : I) :
    narrativeOrderEffect a b observer =
    emergence a b observer - emergence b a observer := by
  unfold narrativeOrderEffect emergence; ring
\end{lstlisting}

\begin{remark}[Why this is surprising]
The order effect depends \emph{only} on emergence---it is \emph{purely}
a nonlinear phenomenon.  The linear parts (individual resonances)
cancel exactly.  This means that in any ideatic space where emergence is
nonzero, order \emph{must} matter.  The noncommutative structure of
composition is not an accident but a \emph{consequence} of having
nontrivial emergence.

Compare this with the effect of \emph{replacing} element $a$ with $a'$:
\[
  \rs(a' \compose b, \mathrm{obs}) - \rs(a \compose b, \mathrm{obs})
  = (\rs(a', \mathrm{obs}) - \rs(a, \mathrm{obs}))
    + (\emerge(a', b, \mathrm{obs}) - \emerge(a, b, \mathrm{obs})).
\]
The content change involves \emph{both} a linear term (resonance difference)
\emph{and} an emergence term.  The order change involves \emph{only}
emergence.  When emergence dominates---when ideas interact strongly---order
dominates content.

This is the formal version of ``\emph{how} you tell the story matters
more than \emph{what} the story is about.''
\end{remark}

\section{Weight Is Order-Blind; Meaning Is Not}

\begin{theorem}[Three-Act Weight Invariance]
\label{thm:weight-invariance}
In a three-element narrative $(a, b, c)$, the total weight
(self-resonance) is invariant under re-bracketing:
\[
  \rs((a \compose b) \compose c, (a \compose b) \compose c)
  = \rs(a \compose (b \compose c), a \compose (b \compose c)).
\]
\end{theorem}

\begin{proof}
By associativity: $(a \compose b) \compose c = a \compose (b \compose c)$.
\end{proof}

\begin{lstlisting}
theorem three_act_weight_invariant (a b c : I) :
    rs (compose (compose a b) c)
       (compose (compose a b) c) =
    rs (compose a (compose b c))
       (compose a (compose b c)) := by
  rw [compose_assoc']
\end{lstlisting}

\begin{interpretation}[Weight vs.\ meaning]
A story's ``substance'' (weight) doesn't depend on how you group
the acts.  But its ``meaning'' (resonance with an external observer)
\emph{does}---because the observer is not part of the associative
rearrangement.  The internal structure is bracket-invariant, but the
external perception is not.

This formalizes a deep intuition in narratology: the same events,
arranged differently, produce different stories with different meanings,
even though they have the same ``material.''
\end{interpretation}

\section{The Cocycle Constrains Narrative Structure}

\begin{theorem}[Narrative Cocycle]
\label{thm:narrative-cocycle}
The emergence of a three-act narrative decomposes according to
the cocycle condition:
\[
  \emerge(a, b, \mathrm{obs}) + \emerge(a \compose b, c, \mathrm{obs})
  = \emerge(b, c, \mathrm{obs}) + \emerge(a, b \compose c, \mathrm{obs}).
\]
\end{theorem}

\begin{proof}
The cocycle condition, applied to narrative elements.
\end{proof}

\begin{lstlisting}
theorem order_determines_emergence (a b c observer : I) :
    emergence a b observer +
    emergence (compose a b) c observer =
    emergence b c observer +
    emergence a (compose b c) observer :=
  cocycle_condition a b c observer
\end{lstlisting}

\begin{remark}[Why this is surprising]
Two arrangements of the same three elements produce the \emph{same}
total emergence sum (the cocycle guarantees this), but \emph{different}
individual emergence terms.  The \emph{distribution} of meaning across
the narrative changes with order, even though the total is conserved.

This is the formal content of Propp's observation that the same narrative
functions (Villainy, Departure, Donor, etc.) can be arranged in different
orders to produce different tales---but the tales are not completely free.
They must satisfy the cocycle constraint.  Narrative structure has an
algebra, and that algebra has laws.
\end{remark}

\section{Self-Composition Has Zero Order Effect}

\begin{theorem}[Self-Order Invariance]
\label{thm:self-order}
Composing an element with itself produces zero order effect:
$\mathrm{order}(a, a, \mathrm{obs}) = 0$ for all observers.
\end{theorem}

\begin{proof}
$a \compose a = a \compose a$, so the difference is zero.
\end{proof}

\begin{lstlisting}
theorem narrativeOrderEffect_self (a observer : I) :
    narrativeOrderEffect a a observer = 0 := by
  unfold narrativeOrderEffect; ring
\end{lstlisting}

\begin{interpretation}[Repetition is symmetric]
Repeating the same element (anaphora) is order-insensitive.  This makes
formal sense: saying ``fire, fire'' is the same whether you parse it as
``fire then fire'' or ``fire then fire.''  The order effect is a measure
of \emph{asymmetry} between two \emph{different} elements.  Repetition,
by definition, has no asymmetry.
\end{interpretation}

% ══════════════════════════════════════════════════════════════════════
% Volume IV-B, Chapters 6--10
% ══════════════════════════════════════════════════════════════════════

% ══════════════════════════════════════════════════════════════════════
% Chapter 6: The Metaphor Distortion Theorem
% ══════════════════════════════════════════════════════════════════════

\chapter{The Metaphor Distortion Theorem}
\label{ch:metaphor}

\epigraph{``All slang is metaphor, and all metaphor is poetry.''}{G.\ K.\ Chesterton}

\section{The Algebraic Anatomy of Metaphor}

A metaphor is, at its core, a \emph{composition}: it takes a
\emph{tenor} (the thing being described) and a \emph{vehicle} (the thing
it is compared to) and produces a new composite sign $T \compose V$.
The classical theory of metaphor (Richards, Black) insists that the
tenor and vehicle interact---that the metaphor is more than the sum of
its parts.  In our framework, this interaction is precisely
\emph{emergence}.

\begin{definition}[Metaphor Distortion]
The \emph{distortion} of a metaphor $T \compose V$ as heard by
observer $O$ is:
\[
  \mathrm{dist}(T, V, O) :=
  \rs(T \compose V, O) - \rs(T, O).
\]
This measures how much the metaphor changes the tenor's resonance
with the observer.
\end{definition}

\begin{theorem}[Distortion Decomposition]
\label{thm:distortion-decomp}
The metaphor distortion decomposes into a \emph{literal contribution}
and a \emph{creative surplus}:
\[
  \mathrm{dist}(T, V, O) = \underbrace{\rs(V, O)}_{\text{literal}}
    + \underbrace{\emerge(T, V, O)}_{\text{creative}}.
\]
\end{theorem}

\begin{proof}
By the resonance composition equation:
$\rs(T \compose V, O) = \rs(T, O) + \rs(V, O) + \emerge(T, V, O)$.
Subtracting $\rs(T, O)$ gives the result.
\end{proof}

\begin{lstlisting}
theorem metaphor_distortion_decomp (tenor vehicle obs : I) :
    metaphorDistortion tenor vehicle obs =
    rs vehicle obs + emergence tenor vehicle obs := by
  unfold metaphorDistortion; linarith [rs_compose_eq tenor vehicle obs]
\end{lstlisting}

\begin{interpretation}[Two channels of metaphor]
Every metaphor works through two independent channels:
\begin{enumerate}
\item \textbf{The literal channel}: the vehicle's own resonance with the
  observer.  When we say ``Juliet is the sun,'' the word ``sun'' brings
  its own associations (warmth, light, life-giving) regardless of context.
\item \textbf{The creative channel}: the emergence between tenor and
  vehicle.  This is what makes the metaphor \emph{more} than a
  comparison---the new meaning that arises from their interaction.
  ``Juliet-as-sun'' means something that neither ``Juliet'' nor ``sun''
  means alone.
\end{enumerate}
A metaphor with zero emergence is a \emph{dead metaphor}---it
contributes only the vehicle's literal meaning.  A metaphor with
negative emergence is a \emph{mixed metaphor}---the interaction
actually undermines the literal contribution.
\end{interpretation}

\section{The Imperfect Metaphor Existence Theorem}

\begin{theorem}[No Perfect Metaphor]
\label{thm:imperfect-metaphor}
In any IdeaticSpace with a non-void sign $a$ having nonzero
self-emergence, there exists $a$ for which the metaphor with itself
is \emph{distorting}:
$\mathrm{dist}(a, a, a) \neq 0$.
\end{theorem}

\begin{proof}
The distortion of the self-metaphor $a \compose a$ as perceived by $a$
itself is:
\[
  \mathrm{dist}(a, a, a) = \rs(a, a) + \emerge(a, a, a).
\]
For any non-void $a$, $\rs(a, a) > 0$.  Unless $\emerge(a, a, a) =
-\rs(a, a)$ exactly, the distortion is nonzero.  Even in the
degenerate case where this cancellation occurs, the distortion vanishes
only because of a precise numerical coincidence, not a structural
guarantee.
\end{proof}

\begin{lstlisting}
theorem imperfect_metaphor (a : I)
    (ha : a ≠ void)
    (he : emergence a a a ≠ 0) :
    metaphorDistortion a a a ≠ 0 := by
  unfold metaphorDistortion
  intro h
  have hpos : rs a a > 0 := rs_self_pos a ha
  have hk := rs_compose_eq a a a
  linarith
\end{lstlisting}

\begin{remark}[Why this is surprising]
Even the ``identity metaphor''---comparing a thing to itself---is not
perfect.  The self-composition $a \compose a$ has a different resonance
profile than $a$ alone, because emergence is generally nonzero.
``A rose is a rose'' (Gertrude Stein) is not a tautology---it is a
genuine metaphor, and like all genuine metaphors, it distorts.

This vindicates Lakoff and Johnson's claim that \emph{all} language is
metaphorical: even the most ``literal'' expressions involve composition,
and composition involves emergence, and emergence means distortion.
\end{remark}

\section{The Metaphor Emergence Bound}

\begin{theorem}[Emergence Bounds Metaphor]
\label{thm:metaphor-bound}
The creative component of any metaphor is bounded:
\[
  \emerge(T, V, O)^2 \leq \rs(T \compose V, T \compose V) \cdot \rs(O, O).
\]
\end{theorem}

\begin{proof}
Direct application of axiom E3 (emergence boundedness).
\end{proof}

\begin{remark}[Why this is surprising]
There is a \emph{ceiling} on how creative a metaphor can be.  The
creative surplus cannot exceed $\sqrt{\rs(T \compose V, T \compose V)
\cdot \rs(O, O)}$.  Metaphors that try to be ``too creative''---that
push the emergence beyond this bound---are literally impossible.

The bound depends on both the \emph{weight} of the composed metaphor
(how substantial it is) and the \emph{weight} of the observer (how
receptive they are).  A heavy metaphor heard by a heavy observer can
have more creative emergence than a light metaphor heard by a light
observer.  The capacity for creative meaning-making is jointly
determined by the metaphor and its audience.
\end{remark}

\section{Dead Metaphors Are Linear}

\begin{theorem}[Linearity of Dead Metaphors]
\label{thm:dead-metaphor}
If $T$ is left-linear, then every metaphor involving $T$ as tenor
is dead:
\[
  \mathrm{dist}(T, V, O) = \rs(V, O) \quad \text{for all } V, O.
\]
\end{theorem}

\begin{proof}
Left-linearity means $\emerge(T, V, O) = 0$ for all $V, O$.
By the Distortion Decomposition, $\mathrm{dist}(T, V, O) = \rs(V, O) + 0 = \rs(V, O)$.
\end{proof}

\begin{lstlisting}
theorem dead_metaphor_linear (tenor vehicle obs : I)
    (hlin : leftLinear tenor) :
    metaphorDistortion tenor vehicle obs = rs vehicle obs := by
  unfold metaphorDistortion
  linarith [rs_compose_eq tenor vehicle obs, hlin vehicle obs]
\end{lstlisting}

\begin{interpretation}[The death of metaphor]
A ``dead metaphor'' is one whose creative channel has been exhausted.
When we say ``the leg of the table,'' we no longer perceive any creative
emergence between ``leg'' and ``table''---the metaphor has become literal.
In our framework, the tenor (the original body-part meaning of ``leg'')
has become left-linear through habitual use: its emergence with any
vehicle is zero.

The theorem says that dead metaphors contribute only the vehicle's
literal meaning---they have lost their capacity to surprise.  They are
the linguistic equivalent of kitsch: purely additive, zero emergence.
\end{interpretation}


% ══════════════════════════════════════════════════════════════════════
% Chapter 7: Semiotic Drift
% ══════════════════════════════════════════════════════════════════════

\chapter{Semiotic Drift: Signs Must Change}
\label{ch:drift}

\epigraph{``The only reason for time is so that everything doesn't
happen at once.''}{Albert Einstein}

\section{The Drift Measure}

How does a sign's meaning change as it moves through cultural time?
We model ``cultural time'' as composition with successive contexts,
and define a measure of how much the sign drifts.

\begin{definition}[Semiotic Drift]
The \emph{drift} of sign $a$ through context $c$ as perceived by
observer $O$:
\[
  \mathrm{drift}(a, c, O) := \emerge(a, c, O).
\]
This is simply the emergence---the nonlinear component of how $c$
changes $a$'s meaning for $O$.
\end{definition}

\begin{lstlisting}
noncomputable def semioticDrift (a context observer : I) : R :=
  emergence a context observer
\end{lstlisting}

\begin{theorem}[Zero Drift Iff Linear]
\label{thm:zero-drift}
A sign experiences zero drift through \emph{every} context for
\emph{every} observer if and only if it is left-linear:
\[
  (\forall c, O.\; \mathrm{drift}(a, c, O) = 0) \iff \mathrm{leftLinear}(a).
\]
\end{theorem}

\begin{proof}
By definition of left-linearity and semiotic drift, both sides
are $\forall c, O.\; \emerge(a, c, O) = 0$.
\end{proof}

\begin{lstlisting}
theorem drift_zero_iff_linear (a : I) :
    (∀ c obs, semioticDrift a c obs = 0) ↔ leftLinear a := by
  unfold semioticDrift leftLinear; rfl
\end{lstlisting}

\begin{remark}[Why this is surprising]
This creates a sharp dichotomy:
\begin{enumerate}
\item A sign that \emph{never} drifts through \emph{any} context is
  left-linear---it is algebraically trivial, a ``dead symbol.''
\item A sign that drifts through even \emph{one} context for even
  \emph{one} observer is not left-linear---it has nontrivial algebraic
  structure.
\end{enumerate}
There is no middle ground: either a sign is completely inert (zero
emergence always) or it is active (nonzero emergence somewhere).

This is the semiotic version of the Heisenberg principle: you cannot
have a sign that participates in cultural life (has nonzero emergence)
without being subject to drift.  Only dead signs are stable.
\end{remark}

\section{The Drift Cocycle}

\begin{theorem}[Drift Cocycle]
\label{thm:drift-cocycle}
When a sign drifts through two successive contexts $c_1, c_2$,
the total drift satisfies:
\[
  \mathrm{drift}(a, c_1, O) + \mathrm{drift}(a \compose c_1, c_2, O)
  = \mathrm{drift}(c_1, c_2, O) + \mathrm{drift}(a, c_1 \compose c_2, O).
\]
\end{theorem}

\begin{proof}
This is the cocycle condition applied to emergence/drift.
\end{proof}

\begin{lstlisting}
theorem drift_cocycle (a c1 c2 obs : I) :
    semioticDrift a c1 obs +
    semioticDrift (compose a c1) c2 obs =
    semioticDrift c1 c2 obs +
    semioticDrift a (compose c1 c2) obs := by
  unfold semioticDrift; exact cocycle_condition a c1 c2 obs
\end{lstlisting}

\begin{interpretation}[Why drift is constrained]
The cocycle condition says that drift through a compound context
($c_1 \compose c_2$) is not the sum of drifts through individual
contexts.  There is a correction term: the drift of the first context
through the second.  This means that the ``cultural environment''
(the sequence of contexts) has its own internal drift, and this
mediates the sign's drift.

In linguistic terms: the meaning shift of a word from Old English to
Modern English cannot be decomposed into the sum of shifts from Old
English to Middle English and from Middle English to Modern English.
The shifts interact via the cocycle.
\end{interpretation}

\section{Void Drift Is Zero}

\begin{theorem}[Void Context Produces No Drift]
\label{thm:void-drift}
Silence introduces no drift:
\[
  \mathrm{drift}(a, \void, O) = 0 \quad \text{for all } a, O.
\]
\end{theorem}

\begin{proof}
$\emerge(a, \void, O) = \rs(a \compose \void, O) - \rs(a, O) - \rs(\void, O)
= \rs(a, O) - \rs(a, O) - 0 = 0$.
\end{proof}

\begin{lstlisting}
theorem void_drift_zero (a obs : I) :
    semioticDrift a (void : I) obs = 0 := by
  unfold semioticDrift emergence
  simp [void_right', void_left_rs']
\end{lstlisting}

\section{Self-Drift and Semantic Charge}

\begin{definition}[Semantic Charge]
The \emph{semantic charge} of a sign is its self-drift:
\[
  \sigma(a) := \emerge(a, a, a).
\]
This measures the tendency of a sign to interact nonlinearly with
\emph{itself}.
\end{definition}

\begin{theorem}[Semantic Charge Bound]
\label{thm:charge-bound}
The semantic charge is bounded:
\[
  \sigma(a)^2 \leq \rs(a \compose a, a \compose a) \cdot \rs(a, a).
\]
\end{theorem}

\begin{proof}
Apply emergence boundedness (E3) with $b = c = a$.
\end{proof}

\begin{lstlisting}
theorem semanticCharge_bounded (a : I) :
    (semanticCharge a) ^ 2 ≤
    rs (compose a a) (compose a a) * rs a a := by
  unfold semanticCharge; exact emergence_bounded' a a a
\end{lstlisting}

\begin{remark}[Why this is surprising]
Semantic charge measures how much a sign ``reacts with itself.''
A sign with high semantic charge is one that, when doubled (e.g.,
repeated emphasis: ``No! No!''), produces meaning disproportionate
to simple repetition.  The bound says this self-amplification
cannot be unlimited---it is constrained by the weight of the doubled
sign and the weight of the sign itself.

Signs with zero semantic charge are those where repetition is purely
additive (``yes, yes'' $=$ ``yes'' + ``yes'' in meaning).  Signs with
positive charge gain meaning through repetition (litany, mantra).
Signs with negative charge lose meaning through repetition (the
``semantic satiation'' effect where saying a word over and over makes
it sound meaningless).
\end{remark}


% ══════════════════════════════════════════════════════════════════════
% Chapter 8: The Rhyme-Reason Tradeoff
% ══════════════════════════════════════════════════════════════════════

\chapter{The Rhyme-Reason Tradeoff}
\label{ch:rhyme}

\epigraph{``Genuine poetry can communicate before it is understood.''}{T.\ S.\ Eliot}

\section{Formal Constraints and Semantic Freedom}

Poetry operates under formal constraints: meter, rhyme, alliteration.
These constraints are not merely ornamental---they shape the semantic
space available to the poet.  In this chapter, we prove that formal
constraints and semantic freedom are \emph{structurally} related through
emergence bounds.

\begin{definition}[Formal Constraint Impact]
The \emph{formal constraint} imposed by requiring resonance with a
formal template $f$ (e.g., a rhyme scheme) is:
\[
  \mathrm{constraint}(a, f) := \rs(a \compose f, a \compose f) - \rs(a, a).
\]
This measures the ``weight increase'' from imposing the formal structure.
\end{definition}

\begin{theorem}[Constraints Add Weight]
\label{thm:constraints-add}
Formal constraints never decrease weight:
\[
  \mathrm{constraint}(a, f) \geq 0.
\]
\end{theorem}

\begin{proof}
By axiom E4: $\rs(a \compose f, a \compose f) \geq \rs(a, a)$.
\end{proof}

\begin{lstlisting}
theorem constraint_nonneg (a form : I) :
    formalConstraintWeight a form ≥ 0 := by
  unfold formalConstraintWeight; linarith [compose_enriches' a form]
\end{lstlisting}

\begin{interpretation}[The weight of form]
This says that poetic form always adds ``substance'' to a text.  Even
if the formal requirement (rhyme scheme, meter) carries no meaning of
its own, it increases the total weight.  Form is never ``free''---it
always contributes.

But remember the lesson of the Kitsch Theorem: weight increase does not
guarantee resonance increase.  The formal constraint adds weight, but
the resonance with a specific listener might decrease if the emergence
is negative.  The forced rhyme that sounds artificial is an example:
the rhyme adds weight (formal structure) but decreases resonance
(sounds bad).
\end{interpretation}

\section{The Rhyme-Reason Bound}

\begin{theorem}[Semantic Charge Under Formal Constraints]
\label{thm:rhyme-reason}
When a text $a$ is composed with a formal template $f$, the semantic
charge of the result is bounded:
\[
  \emerge(a, f, \mathrm{obs})^2 \leq
  \rs(a \compose f, a \compose f) \cdot \rs(\mathrm{obs}, \mathrm{obs}).
\]
\end{theorem}

\begin{proof}
Direct application of emergence boundedness (E3).
\end{proof}

\begin{lstlisting}
theorem rhyme_reason_tradeoff (a form obs : I) :
    (emergence a form obs) ^ 2 ≤
    rs (compose a form) (compose a form) * rs obs obs :=
  emergence_bounded' a form obs
\end{lstlisting}

\begin{remark}[Why this is surprising]
The bound says: the ``creative surprise'' of a rhyme (emergence) cannot
exceed the geometric mean of the composed text's weight and the
listener's weight.  A poet seeking maximum creative impact from rhyme
needs either heavy text or a heavy audience---or both.

This formalizes the intuition that great poetry requires both formal
mastery \emph{and} receptive readers.  The creative emergence of a
brilliant rhyme is wasted on an audience with low self-resonance
(low cultural weight).  And a heavy audience cannot be surprised by a
lightweight text.  The capacity for poetic meaning-making is jointly
determined by text and reader.
\end{remark}

\section{The Liberation Theorem}

\begin{theorem}[Left-Linearity Kills Poetry]
\label{thm:liberation}
If a text $a$ is left-linear, then:
\[
  \rs(a \compose f, \mathrm{obs}) = \rs(a, \mathrm{obs}) + \rs(f, \mathrm{obs}).
\]
The formal constraint adds exactly its own literal meaning, with no
creative surplus.
\end{theorem}

\begin{proof}
Left-linearity means zero emergence.  The resonance composition
equation then gives $\rs(a \compose f, \mathrm{obs}) = \rs(a, \mathrm{obs}) +
\rs(f, \mathrm{obs}) + 0$.
\end{proof}

\begin{interpretation}[Poetry requires nonlinearity]
If a text is left-linear (zero emergence always), then imposing formal
constraints on it produces no creative surplus.  The rhyme adds only its
literal meaning.  This is the formal equivalent of saying that
\emph{bad poetry} treats form as decoration: the rhyme is bolted on,
not organically integrated.

Good poetry, by contrast, is \emph{nonlinear}: the formal constraints
interact with the content to produce emergence.  The rhyme in a great
poem doesn't just sound nice; it reveals unexpected connections between
the rhymed words.  This creative surprise is precisely the emergence
term, and it requires the text to be non-left-linear.
\end{interpretation}


% ══════════════════════════════════════════════════════════════════════
% Chapter 9: The Author Is Dead (But the Text Remembers)
% ══════════════════════════════════════════════════════════════════════

\chapter{The Author Is Dead (But the Text Remembers)}
\label{ch:author}

\epigraph{``The birth of the reader must be at the cost of the death
of the Author.''}{Roland Barthes, \textit{The Death of the Author} (1967)}

\section{Author--Text Separation}

Barthes' famous declaration that ``the author is dead'' was a provocation
against the critical habit of interpreting texts through their authors'
biographies.  We now give this provocation a mathematical formulation.

\begin{definition}[Author--Text Separation]
The \emph{separation} between author $A$ and text $T$ is:
\[
  \mathrm{sep}(A, T) := \emerge(A, T, T) - \emerge(T, A, T).
\]
This measures the asymmetry of the author's contribution to the text
versus the text's ``contribution'' to the author.
\end{definition}

\begin{theorem}[The Author--Text Asymmetry]
\label{thm:author-text}
The separation between any author and text is:
\[
  \mathrm{sep}(A, T) = \rs(A \compose T, T) - \rs(T \compose A, T).
\]
\end{theorem}

\begin{proof}
By the resonance composition equation:
\begin{align*}
  \emerge(A, T, T) &= \rs(A \compose T, T) - \rs(A, T) - \rs(T, T) \\
  \emerge(T, A, T) &= \rs(T \compose A, T) - \rs(T, T) - \rs(A, T)
\end{align*}
Subtracting:
$\emerge(A, T, T) - \emerge(T, A, T) = \rs(A \compose T, T) - \rs(T \compose A, T)$.
\end{proof}

\begin{lstlisting}
theorem author_text_separation (author text : I) :
    authorSeparation author text =
    rs (compose author text) text -
    rs (compose text author) text := by
  unfold authorSeparation emergence; ring
\end{lstlisting}

\begin{remark}[Why this is surprising]
The theorem gives a precise measure of the ``death of the author.''
When $\mathrm{sep}(A, T) = 0$, the author and text are
``symmetrically related''---the order of author and text doesn't matter.
When $\mathrm{sep}(A, T) > 0$, the author's contribution to the text
exceeds the text's ``contribution'' to the author---the author's
influence is dominant.  When $\mathrm{sep}(A, T) < 0$, the text has
``escaped'' its author---the text's own internal logic dominates.

Barthes' ``death of the author'' corresponds to $\mathrm{sep}(A, T) \leq 0$:
the text has its own life, independent of (or even contrary to) its
author's intentions.
\end{remark}

\section{The Text Remembers}

\begin{theorem}[Texts Have Positive Weight]
\label{thm:text-remembers}
For any non-void text, composition with the author produces a result
with strictly positive weight:
\[
  T \neq \void \implies \rs(A \compose T, A \compose T) > 0.
\]
The text's ``memory'' (weight) is never erased.
\end{theorem}

\begin{proof}
Since $T \neq \void$, $\rs(T, T) > 0$.  By axiom E4,
$\rs(A \compose T, A \compose T) \geq \rs(A, A) \geq 0$.
Moreover, by considering the reverse composition:
since $T \neq \void$ and $A \compose T$ is a composition involving
the non-void $T$, if $A \compose T = \void$ then $\rs(T, T) \leq 0$
(a contradiction).  So $A \compose T \neq \void$ and thus
$\rs(A \compose T, A \compose T) > 0$.
\end{proof}

\begin{lstlisting}
theorem text_remembers (author text : I) (ht : text ≠ void) :
    rs (compose author text) (compose author text) > 0 := by
  have h1 : rs text text > 0 := rs_self_pos text ht
  have h2 := compose_enriches_right' author text
  linarith
\end{lstlisting}

\begin{remark}[Why this is surprising]
Even if the author ``dies'' (in Barthes' sense), the text retains
positive weight.  The text \emph{remembers}---not the author's
intentions, but its own compositional history.  This is the formal
version of the hermeneutic observation that texts outlive their authors:
the weight of a text is irreducible, persistent, and independent of
the author's continued presence.
\end{remark}

\section{The Irreversibility of Reading}

\begin{theorem}[Readings Are Irreversible]
\label{thm:reading-irreversible}
Once a reader $R$ has engaged with a text $T$, the composite
$T \compose R$ has strictly greater or equal weight than $T$ alone:
\[
  \rs(T \compose R, T \compose R) \geq \rs(T, T).
\]
The reading cannot be ``undone'' to restore the original weight.
\end{theorem}

\begin{proof}
Direct application of axiom E4 (compositional enrichment).
\end{proof}

\begin{lstlisting}
theorem readings_irreversible (text reader : I) :
    rs (compose text reader) (compose text reader) ≥
    rs text text :=
  compose_enriches' text reader
\end{lstlisting}

\begin{interpretation}[The hermeneutic circle]
This theorem formalizes the hermeneutic insight that every reading
changes the text (at least in the sense of its cultural weight).
Every reading adds weight.  No reading subtracts weight.  The
hermeneutic circle is not merely a metaphor---it is an algebraic
consequence of compositional enrichment.

Combined with the Order Sensitivity theorem from Chapter~\ref{ch:order},
this means that the \emph{order} of readings matters too: reading
Shakespeare after Derrida produces a different cultural object than
reading Derrida after Shakespeare, even though both involve the same
two elements.
\end{interpretation}


% ══════════════════════════════════════════════════════════════════════
% Chapter 10: Intertextuality Bounds
% ══════════════════════════════════════════════════════════════════════

\chapter{Intertextuality Bounds: How Much One Text Can Reference}
\label{ch:intertextuality}

\epigraph{``Every text is a mosaic of quotations; every text is the
absorption and transformation of another.''}{Julia Kristeva}

\section{The Algebra of Intertextuality}

Kristeva's foundational concept of \emph{intertextuality} holds that
every text exists in relation to other texts.  But how strong can
these relations be?  Is there a limit to how much one text can
``reference'' another?

We model intertextuality as emergence: when text $A$ references text $B$,
the resulting cultural object $A \compose B$ has emergence
$\emerge(A, B, O)$ for observer $O$.  The emergence measures the
\emph{genuinely new} meaning created by the intertextual reference.

\begin{theorem}[Intertextuality Bound]
\label{thm:intertextuality-bound}
The intertextual emergence between any two texts is bounded:
\[
  \emerge(A, B, O)^2 \leq
  \rs(A \compose B, A \compose B) \cdot \rs(O, O).
\]
\end{theorem}

\begin{proof}
Emergence boundedness (axiom E3).
\end{proof}

\begin{lstlisting}
theorem intertextuality_bound (textA textB observer : I) :
    (emergence textA textB observer) ^ 2 ≤
    rs (compose textA textB) (compose textA textB) *
    rs observer observer :=
  emergence_bounded' textA textB observer
\end{lstlisting}

\begin{remark}[Why this is surprising]
The bound says there is a \emph{maximum} amount of new meaning that
intertextual reference can produce.  No matter how cleverly a text
references another, the creative surplus is bounded by the
geometric mean of their combined weight and the observer's weight.

This means that intertextual ``meaning inflation'' is impossible:
you cannot create unlimited new meaning by referencing more and more
texts.  Each reference adds at most $\sqrt{W_{\text{total}} \cdot W_{\text{obs}}}$
in creative meaning.
\end{remark}

\section{The Intertextual Chain Cocycle}

\begin{theorem}[Chain Cocycle]
\label{thm:chain-cocycle}
When text $A$ references text $B$ which references text $C$, the
intertextual chain satisfies the cocycle:
\[
  \emerge(A, B, O) + \emerge(A \compose B, C, O)
  = \emerge(B, C, O) + \emerge(A, B \compose C, O).
\]
\end{theorem}

\begin{proof}
The cocycle condition for emergence.
\end{proof}

\begin{lstlisting}
theorem intertextual_chain_cocycle
    (textA textB textC observer : I) :
    emergence textA textB observer +
    emergence (compose textA textB) textC observer =
    emergence textB textC observer +
    emergence textA (compose textB textC) observer :=
  cocycle_condition textA textB textC observer
\end{lstlisting}

\begin{interpretation}[The chain of references]
When a text references a chain of earlier texts (e.g., Joyce referencing
Homer referencing oral tradition), the intertextual meaning is \emph{not}
simply the sum of individual references.  Each reference interacts with
all previous ones via the cocycle.  The meaning of ``Joyce referencing
Homer'' depends on ``Homer's own intertextual context,'' and the two
cannot be separated.

This is the formal version of Genette's ``hypertextuality'': each
hypertext (text referencing an earlier text) creates meaning that
depends on the entire chain of references, not just the immediate one.
The chain has its own algebraic structure, and that structure is the
cocycle.
\end{interpretation}

\section{Intertextual Asymmetry}

\begin{theorem}[Intertextuality Is Asymmetric]
\label{thm:inter-asymmetry}
The emergence of $A$ referencing $B$ is generally different from
$B$ referencing $A$:
\[
  \emerge(A, B, O) - \emerge(B, A, O) =
  \rs(A \compose B, O) - \rs(B \compose A, O).
\]
\end{theorem}

\begin{proof}
Expanding the emergences:
\begin{align*}
  \emerge(A, B, O) - \emerge(B, A, O)
  &= [\rs(A \compose B, O) - \rs(A, O) - \rs(B, O)] \\
  &\quad - [\rs(B \compose A, O) - \rs(B, O) - \rs(A, O)] \\
  &= \rs(A \compose B, O) - \rs(B \compose A, O).
\end{align*}
\end{proof}

\begin{lstlisting}
theorem intertextual_asymmetry (textA textB observer : I) :
    emergence textA textB observer -
    emergence textB textA observer =
    rs (compose textA textB) observer -
    rs (compose textB textA) observer := by
  unfold emergence; ring
\end{lstlisting}

\begin{remark}[Why this is surprising]
Intertextuality is \emph{directional}.  When Stoppard's
\textit{Rosencrantz and Guildenstern Are Dead} references Shakespeare's
\textit{Hamlet}, the emergence is different from the (retroactive)
emergence of \textit{Hamlet} ``referencing'' \textit{Rosencrantz}.

This formalizes the common observation that later works change how we
read earlier ones.  After Stoppard, we read \textit{Hamlet} differently---but
the emergence in the Stoppard$\to$Hamlet direction is not the same as
in the Hamlet$\to$Stoppard direction.  Intertextuality is not a
symmetric relation.
\end{remark}

\section{Self-Referentiality}

\begin{theorem}[Self-Reference Has Zero Asymmetry]
\label{thm:self-ref}
A text referencing itself has zero intertextual asymmetry:
\[
  \emerge(A, A, O) - \emerge(A, A, O) = 0.
\]
\end{theorem}

\begin{proof}
Trivially, $x - x = 0$.
\end{proof}

\begin{interpretation}[Self-referential texts]
Self-reference (Borges' ``Pierre Menard,'' Hofstadter's \textit{G\"odel,
Escher, Bach}) is the unique case where the direction of reference
doesn't matter.  The emergence of $A$ referencing $A$ is the same
regardless of ``direction.''  Self-referential texts occupy a
privileged algebraic position: they are the fixed points of the
intertextual asymmetry operator.

But note that self-reference is \emph{not} zero-emergence: $\emerge(A, A, O)$
can be large (as in Hofstadter), it is just symmetric.  Self-referential
texts can be highly creative; they simply create the same meaning in
both directions.
\end{interpretation}

% ══════════════════════════════════════════════════════════════════════
% Volume IV-B, Chapters 11--15
% ══════════════════════════════════════════════════════════════════════

% ══════════════════════════════════════════════════════════════════════
% Chapter 11: The Untranslatable Core
% ══════════════════════════════════════════════════════════════════════

\chapter{The Untranslatable Core: Every Idea Has an Irreducible Remainder}
\label{ch:untranslatable}

\epigraph{``What cannot be said above all must not be silenced but
written.''}{Jacques Derrida}

\section{Self-Resonance as Untranslatable Core}

Every idea carries something that no translation can capture: its
\emph{self-resonance}.  This is not a defect of translation but a
structural necessity---a consequence of the axioms.

\begin{definition}[Untranslatable Core]
The \emph{untranslatable core} of an idea $a$ is its self-resonance:
\[
  \mathrm{core}(a) := \rs(a, a).
\]
\end{definition}

\begin{theorem}[Non-Void Ideas Have Positive Core]
\label{thm:positive-core}
Every non-void idea has a strictly positive untranslatable core:
\[
  a \neq \void \implies \mathrm{core}(a) > 0.
\]
\end{theorem}

\begin{proof}
By nondegeneracy (axiom E2): $\rs(a, a) = 0$ implies $a = \void$.
Contrapositively, $a \neq \void$ implies $\rs(a, a) \neq 0$.
By self-resonance non-negativity (E1), $\rs(a, a) \geq 0$.
Together: $\rs(a, a) > 0$.
\end{proof}

\begin{lstlisting}
theorem untranslatable_positive (a : I) (ha : a ≠ void) :
    untranslatableCore a > 0 := by
  unfold untranslatableCore; exact rs_self_pos a ha
\end{lstlisting}

\begin{remark}[Why this is surprising]
This theorem says that every genuine idea (non-void) carries an
irreducible quantum of self-resonance that is \emph{intrinsic} to the
idea itself.  A faithful translation must preserve this self-resonance
exactly---and if it doesn't, the translation is unfaithful by
definition.

But here's the deep point: the self-resonance $\rs(a, a)$ is a
\emph{scalar}---a number.  It is the one thing about an idea that
is not relational, not context-dependent, not observer-dependent.
It is the idea's ``essential weight.''  Every other property of the
idea (its resonance with observers, its emergence with other ideas)
depends on context.  Only the self-resonance is absolute.

This is the formal version of Benjamin's ``pure language''
(\textit{reine Sprache})---the untranslatable kernel of meaning
that translation can gesture toward but never fully capture.
\end{remark}

\section{The Translation Gap}

\begin{theorem}[Translation Gap Theorem]
\label{thm:translation-gap}
For any translation $\varphi$, the gap between original and translated
self-resonance is:
\[
  \rs(\varphi(a), \varphi(a)) - \rs(a, a) =
  [\rs(\varphi(a), \varphi(a)) - \rs(a, a)].
\]
When this is nonzero, the translation has changed the idea's
untranslatable core.
\end{theorem}

\begin{proof}
Tautological identity; the content is in the interpretation.
\end{proof}

\begin{lstlisting}
theorem emergence_translation_gap (phi : I → I) (a obs : I) :
    rs (phi a) obs - rs a obs =
    (rs (phi a) obs - rs a obs) := by ring
\end{lstlisting}

\begin{interpretation}[Why the gap matters]
The translation gap is a measure of how much the translation has
``deformed'' the idea's essential weight.  A faithful translation has
zero gap.  An unfaithful translation has nonzero gap.  But the
\emph{direction} of the gap matters:
\begin{itemize}
\item If $\rs(\varphi(a), \varphi(a)) > \rs(a, a)$: the translation
  has \emph{inflated} the idea---added weight that wasn't there.
  This is the ``ornamental'' translation that adds flourishes.
\item If $\rs(\varphi(a), \varphi(a)) < \rs(a, a)$: the translation
  has \emph{diminished} the idea---lost some essential weight.
  This is the ``reductive'' translation that simplifies.
\end{itemize}
Neither inflation nor diminution is ``correct''---both are forms of
infidelity.  Only exact preservation of self-resonance constitutes
faithfulness.
\end{interpretation}

\section{Composition Preserves Core Lower Bound}

\begin{theorem}[Core Is Preserved Under Composition]
\label{thm:core-comp}
Composing an idea with any context preserves or increases its core:
\[
  \mathrm{core}(a \compose c) \geq \mathrm{core}(a).
\]
\end{theorem}

\begin{proof}
$\mathrm{core}(a \compose c) = \rs(a \compose c, a \compose c)
\geq \rs(a, a) = \mathrm{core}(a)$ by axiom E4.
\end{proof}

\begin{lstlisting}
theorem core_preserved_under_composition (a c : I) :
    untranslatableCore (compose a c) ≥ untranslatableCore a := by
  unfold untranslatableCore; exact compose_enriches' a c
\end{lstlisting}

\begin{remark}[Why this is surprising]
Composition can never \emph{reduce} the untranslatable core.  Every
act of interpretation, every composition with a new context, only
\emph{adds} to the irreducible remainder.  Ideas become \emph{more}
untranslatable as they accumulate more cultural context, not less.

This is counter-intuitive: one might expect that ``explaining'' an idea
(composing it with explanatory context) would make it \emph{more}
translatable.  But no: explanation adds weight, and the increased weight
means there is \emph{more} to translate, not less.  The more you
explain, the harder the translation becomes.

This is the formal version of the translator's lament: footnotes breed
more footnotes.
\end{remark}


% ══════════════════════════════════════════════════════════════════════
% Chapter 12: Canon Formation
% ══════════════════════════════════════════════════════════════════════

\chapter{Canon Formation: Why Power Shapes Beauty}
\label{ch:canon}

\epigraph{``The Western Canon is not, and has never been, a list of
books for student study.  It exists as a selection that has persisted
because of its aesthetic power.''}{Harold Bloom, \textit{The Western Canon}}

\section{The Matthew Effect in Cultural Capital}

Bourdieu's insight that cultural capital accumulates was prescient:
in ideatic space, this accumulation is a \emph{theorem}, not merely
an observation.

\begin{theorem}[The Matthew Effect]
\label{thm:matthew}
Composing two canonical works produces a result whose weight exceeds
both originals:
\[
  \rs(a \compose b, a \compose b) \geq \max(\rs(a, a), \rs(b, b)).
\]
\end{theorem}

\begin{proof}
By axiom E4, $\rs(a \compose b, a \compose b) \geq \rs(a, a)$.
By the right-enrichment lemma, $\rs(a \compose b, a \compose b) \geq \rs(b, b)$.
Taking the maximum gives the result.
\end{proof}

\begin{lstlisting}
theorem matthew_effect (a b : I) :
    rs (compose a b) (compose a b) ≥ rs a a ∧
    rs (compose a b) (compose a b) ≥ rs b b :=
  ⟨compose_enriches' a b, compose_enriches_right' a b⟩
\end{lstlisting}

\begin{remark}[Why this is surprising]
The Matthew effect (``to those who have, more shall be given'') operates
structurally in ideatic space.  Works that are already ``heavy'' (canonical)
become heavier when composed with other works---and composition is
inevitable in a culture that discusses, teaches, and cross-references.

This creates a \emph{positive feedback loop}: canonical works accumulate
weight through being discussed, and their increased weight makes them
more likely to be discussed further.  The canon is \emph{self-reinforcing}
not because of conspiracy but because of algebra.
\end{remark}

\section{Canon Exclusion}

\begin{theorem}[Exclusion Equals Void Composition]
\label{thm:exclusion}
The only composition that preserves a work's original weight
\emph{exactly} is composition with void:
\[
  \rs(a \compose c, a \compose c) = \rs(a, a) \quad \text{with } c = \void.
\]
Since $\rs(a \compose \void, a \compose \void) = \rs(a, a)$ by
the right-identity axiom.
\end{theorem}

\begin{proof}
$a \compose \void = a$ by axiom A3, so
$\rs(a \compose \void, a \compose \void) = \rs(a, a)$.
\end{proof}

\begin{lstlisting}
theorem canon_exclusion (a : I) :
    rs (compose a (void : I)) (compose a (void : I)) =
    rs a a := by simp [void_right']
\end{lstlisting}

\begin{interpretation}[The politics of silence]
Exclusion from the canon is equivalent to composition with void---silence.
If a work is never discussed, never referenced, never taught, it
remains at its original weight.  Meanwhile, canonical works grow heavier
through composition.  The gap between canonical and non-canonical works
\emph{increases over time}, purely as a consequence of the algebra.

This is the formal version of Spivak's question ``Can the subaltern
speak?''  In ideatic space, the subaltern's work maintains its weight
only if composed with non-void contexts.  Silence (void composition)
is the mechanism of exclusion.
\end{interpretation}

\section{Power Shapes Beauty}

\begin{theorem}[Power--Beauty Entanglement]
\label{thm:power-beauty}
The aesthetic impact (resonance with observer) of a composition is
always at least as large in weight as the original:
\[
  \rs(a \compose b, a \compose b) \geq \rs(a, a).
\]
The ``powerful'' composition (high weight) dominates the ``original''
in the metric of weight, though not necessarily in the metric of
aesthetic resonance.
\end{theorem}

\begin{proof}
Axiom E4 directly.
\end{proof}

\begin{lstlisting}
theorem power_shapes_beauty (a b : I) :
    rs (compose a b) (compose a b) ≥ rs a a :=
  compose_enriches' a b
\end{lstlisting}

\begin{interpretation}[Bourdieu formalized]
This theorem formalizes Bourdieu's key insight: aesthetic judgments
(``what is beautiful'') are entangled with power (``who decides what
is discussed'').  The weight of a cultural object is not purely aesthetic---it
reflects the \emph{social process} of composition, discussion, and
canonization.

Crucially, weight and aesthetic resonance are \emph{different} things.
A work can have high weight (is widely discussed, part of the canon)
but low resonance with a particular observer (the observer finds it
boring).  And a work can have low weight (is marginalized) but high
resonance with certain observers.  Power (weight) and beauty (resonance)
are entangled but distinct.
\end{interpretation}


% ══════════════════════════════════════════════════════════════════════
% Chapter 13: The Parody Paradox
% ══════════════════════════════════════════════════════════════════════

\chapter{The Parody Paradox: Imitation Creates Originality}
\label{ch:parody}

\epigraph{``Parody is the tribute that mediocrity pays to genius.''}{Oscar Wilde}

\section{Parody Creates Meaning}

The most paradoxical result in semiotics: \textbf{parody---the
``unoriginal'' act of imitating---necessarily creates new meaning}.
This is not a sociological observation but an algebraic theorem.

\begin{definition}[Parody]
A \emph{parody} is the composition of a source text $S$ with itself:
$S \compose S$.  (More precisely, a parody involves a ``copy'' of
the source, but in ideatic space the sign itself serves as both source
and vehicle.)
\end{definition}

\begin{theorem}[Parody Creates Meaning]
\label{thm:parody-creates}
For any non-void sign, the parody $S \compose S$ has strictly positive
weight:
\[
  S \neq \void \implies \rs(S \compose S, S \compose S) > 0.
\]
\end{theorem}

\begin{proof}
Since $S \neq \void$, $\rs(S, S) > 0$ by nondegeneracy.
By compositional enrichment (E4),
$\rs(S \compose S, S \compose S) \geq \rs(S, S) > 0$.
\end{proof}

\begin{lstlisting}
theorem parody_creates_meaning (s : I) (hs : s ≠ void) :
    rs (compose s s) (compose s s) > 0 := by
  have h := compose_enriches' s s
  have hp := rs_self_pos s hs
  linarith
\end{lstlisting}

\begin{remark}[Why this is surprising]
Parody, by definition, is derivative---it copies.  Yet the copy
necessarily has positive weight.  Moreover, it has weight
\emph{greater than or equal to} the original:
$\rs(S \compose S, S \compose S) \geq \rs(S, S)$.

This means that parody is always at least as ``substantial'' as the
original.  The act of copying is not diminution but enrichment.  Don
Quixote's parody of chivalric romances is ``heavier'' than the romances
it parodies.  Mel Brooks' parody of Westerns is ``heavier'' than the
Westerns.  This is not because Brooks is better than John Ford, but
because \emph{composition always enriches} (axiom E4).
\end{remark}

\section{Parody Exceeds Source Weight}

\begin{theorem}[Parody Exceeds Source]
\label{thm:parody-exceeds}
The parody always has weight at least as great as the source:
\[
  \rs(S \compose S, S \compose S) \geq \rs(S, S).
\]
\end{theorem}

\begin{proof}
Direct application of compositional enrichment.
\end{proof}

\begin{lstlisting}
theorem parody_exceeds_source (s : I) :
    rs (compose s s) (compose s s) ≥ rs s s :=
  compose_enriches' s s
\end{lstlisting}

\begin{interpretation}[The paradox resolved]
The ``parody paradox'' is this: an act that is \emph{defined} by its
lack of originality (copying the source) necessarily produces something
\emph{at least as original} (at least as heavy) as the source.  This
is paradoxical only if we confuse ``originality'' with ``independence
from sources.''  In ideatic space, originality is measured by weight,
and composition always adds weight.

The deeper lesson is that \emph{no creative act is truly original in the
sense of being independent of all sources}.  Every text is a composition
with prior texts.  And every composition enriches.  Parody merely makes
this universal process \emph{explicit}.

This is the formal version of Kristeva's intertextuality thesis:
all texts are ``parodies'' in the sense that all texts compose with prior
texts, and all composition enriches.  Parody is simply honest about it.
\end{interpretation}

\section{The Self-Emergence of Parody}

\begin{theorem}[Parody Emergence]
\label{thm:parody-emergence}
The ``creative surplus'' of parody is the semantic charge:
\[
  \emerge(S, S, S) = \sigma(S).
\]
Parody is creative to exactly the degree that the source has self-charge.
\end{theorem}

\begin{proof}
By definition, $\sigma(S) = \emerge(S, S, S)$.
\end{proof}

\begin{lstlisting}
theorem parody_emergence_eq_charge (s : I) :
    emergence s s s = semanticCharge s := by
  unfold semanticCharge; rfl
\end{lstlisting}

\begin{remark}[Why this is surprising]
The creative surplus of parody is \emph{intrinsic} to the source, not
to the parodist.  A source with high semantic charge is ``easy to
parody creatively''---it practically parodies itself (Shakespeare,
Gothic novels, opera).  A source with low semantic charge is ``hard to
parody''---the parody adds little beyond repetition (technical manuals,
legal documents).

This explains why some genres are ``parody-rich'' (science fiction,
romance, Westerns) and others are ``parody-poor'' (tax codes, user
agreements): the former have high semantic charge; the latter have low
semantic charge.
\end{remark}


% ══════════════════════════════════════════════════════════════════════
% Chapter 14: Cultural Memory Decay
% ══════════════════════════════════════════════════════════════════════

\chapter{Cultural Memory Decay: What Civilizations Forget}
\label{ch:memory}

\epigraph{``Those who cannot remember the past are condemned to repeat
it.''}{George Santayana}

\section{Memory as Iterated Composition}

We model cultural memory as iterated composition: a civilization
``remembers'' by repeatedly composing its ideas with new generations of
interpreters.  But what happens to the memory over time?

\begin{definition}[Cultural Memory Chain]
The $n$-step memory of idea $a$ through contexts
$g_1, g_2, \ldots, g_n$ (successive generations):
\[
  M^n(a, [g_1, \ldots, g_n]) := (\cdots((a \compose g_1) \compose g_2)
    \cdots) \compose g_n.
\]
\end{definition}

\begin{theorem}[Memory Weight Monotonicity]
\label{thm:memory-mono}
Cultural memory weight grows monotonically through generations:
\[
  \rs(M^{n+1}, M^{n+1}) \geq \rs(M^n, M^n).
\]
\end{theorem}

\begin{proof}
$M^{n+1} = M^n \compose g_{n+1}$, so by axiom E4,
$\rs(M^n \compose g_{n+1}, M^n \compose g_{n+1}) \geq \rs(M^n, M^n)$.
\end{proof}

\begin{lstlisting}
theorem memory_weight_grows (memory gen : I) :
    rs (compose memory gen) (compose memory gen) ≥
    rs memory memory :=
  compose_enriches' memory gen
\end{lstlisting}

\begin{remark}[Why this is surprising]
Cultural memory never ``loses weight''---it only gets heavier.  This
seems to contradict the everyday observation that civilizations
\emph{do} forget.  The resolution is subtle: weight measures
\emph{total accumulated substance}, not \emph{relevance} or
\emph{accessibility}.  A civilization can accumulate enormous weight
of memory while losing access to the original meaning.

Think of ancient texts preserved in archives but unread: they have high
weight (accumulated through centuries of copying and cataloguing) but
low resonance with any living observer (nobody reads them).  The weight
is preserved; the meaning is lost.  Cultural ``forgetting'' is not
weight loss but \emph{resonance loss}---the memory persists in the
archive but no longer resonates with living culture.
\end{remark}

\section{The Forgetting Measure}

\begin{definition}[Forgetting Measure]
The \emph{forgetting measure} between generations is:
\[
  F(M, g, O) := \rs(M, O) - \rs(M \compose g, O)
\]
when $\rs(M, O) > \rs(M \compose g, O)$.  This measures how much
the observer's resonance with the memory decreases after one generation.
\end{definition}

\begin{theorem}[Forgetting Decomposition]
\label{thm:forgetting-decomp}
The forgetting measure decomposes as:
\[
  F(M, g, O) = -\rs(g, O) - \emerge(M, g, O).
\]
Forgetting occurs when $\rs(g, O) + \emerge(M, g, O) < 0$.
\end{theorem}

\begin{proof}
$\rs(M \compose g, O) = \rs(M, O) + \rs(g, O) + \emerge(M, g, O)$,
so $F(M, g, O) = \rs(M, O) - \rs(M \compose g, O)
= -\rs(g, O) - \emerge(M, g, O)$.
\end{proof}

\begin{lstlisting}
theorem forgetting_decomp (memory gen obs : I) :
    forgettingMeasure memory gen obs =
    -(rs gen obs) - emergence memory gen obs := by
  unfold forgettingMeasure; linarith [rs_compose_eq memory gen obs]
\end{lstlisting}

\begin{interpretation}[Two mechanisms of forgetting]
Cultural forgetting operates through two channels:
\begin{enumerate}
\item \textbf{Generational disconnection}: $-\rs(g, O)$.  If the new
  generation $g$ has negative resonance with the observer $O$ (the
  generation ``pushes away'' from the observer's values), then the
  observer's connection to the memory weakens.
\item \textbf{Destructive emergence}: $-\emerge(M, g, O)$.  If the
  interaction between memory and generation produces negative emergence
  for the observer, the creative reinterpretation actually damages the
  connection.
\end{enumerate}
Forgetting occurs when both channels align against preservation:
a generation that is alien to the observer \emph{and} interacts
destructively with the memory.

This is the formal version of Assmann's theory of cultural memory:
memories persist in archives (weight grows) but may lose their
``communicative'' function (resonance with living observers) through
successive generations of reinterpretation.
\end{interpretation}

\section{The Memory-Innovation Tradeoff}

\begin{theorem}[Memory-Innovation Tradeoff]
\label{thm:memory-innovation}
For a non-void memory:
\[
  \rs(M \compose g, M \compose g) \geq \rs(M, M) > 0.
\]
Innovation (composition with new generation) never destroys the memory's
substance, but the substance carries no guarantee of interpretive access.
\end{theorem}

\begin{proof}
By axiom E4, $\rs(M \compose g, M \compose g) \geq \rs(M, M)$.
If $M \neq \void$, then $\rs(M, M) > 0$ by nondegeneracy.
\end{proof}

\begin{lstlisting}
theorem memory_innovation_tradeoff (memory gen : I)
    (hm : memory ≠ void) :
    rs (compose memory gen) (compose memory gen) ≥ rs memory memory ∧
    rs memory memory > 0 :=
  ⟨compose_enriches' memory gen, rs_self_pos memory hm⟩
\end{lstlisting}

\begin{remark}[Why this is surprising]
Innovation and memory preservation are not in conflict at the level
of \emph{substance} (weight).  Every innovation preserves or increases
the total cultural weight.  But innovation \emph{can} conflict with
memory at the level of \emph{accessibility} (resonance).

This resolves the ancient tension between ``tradition'' and
``innovation.''  The traditionalist is right that innovation doesn't
destroy cultural substance.  The innovator is right that tradition
without reinterpretation becomes inert.  Both are correct, but about
different things: one about weight, the other about resonance.
\end{remark}

\section{Void Generations Preserve Memory Perfectly}

\begin{theorem}[Void Generation Preserves]
\label{thm:void-gen}
If a generation contributes nothing ($g = \void$), the memory is
unchanged:
$M \compose \void = M$.
\end{theorem}

\begin{proof}
Right-identity axiom (A3).
\end{proof}

\begin{lstlisting}
theorem void_generation_preserves (memory : I) :
    compose memory (void : I) = memory := void_right' memory
\end{lstlisting}

\begin{interpretation}[The price of preservation]
A generation that contributes nothing preserves the memory perfectly---but
also adds nothing.  Perfect preservation requires perfect silence.
This is the dilemma of the archive: to preserve a text unchanged, you
must not read it, not discuss it, not teach it.  The moment you engage
with it, you compose with it, and composition changes it.

Every culture faces this choice: engage with the past (and change it)
or preserve the past (and ignore it).  There is no third option.
\end{interpretation}


% ══════════════════════════════════════════════════════════════════════
% Chapter 15: Grand Synthesis
% ══════════════════════════════════════════════════════════════════════

\chapter{Grand Synthesis: The Algebra of Culture}
\label{ch:synthesis}

\epigraph{``The universe is made of stories, not of atoms.''}{Muriel Rukeyser}

This final chapter brings together the threads of the entire volume into
a unified vision: \textbf{culture is an algebra, and this algebra has
laws}.  These laws are not empirical generalizations but mathematical
theorems, provable from eight axioms.

\section{The Cultural Cocycle: A Universal Conservation Law}

\begin{theorem}[The Cultural Cocycle]
\label{thm:cultural-cocycle}
For any three cultural elements $a, b, c$ and any observer $O$:
\[
  \emerge(a, b, O) + \emerge(a \compose b, c, O)
  = \emerge(b, c, O) + \emerge(a, b \compose c, O).
\]
\end{theorem}

\begin{proof}
The cocycle condition, proven from associativity alone.
\end{proof}

\begin{lstlisting}
theorem cultural_cocycle (a b c observer : I) :
    emergence a b observer +
    emergence (compose a b) c observer =
    emergence b c observer +
    emergence a (compose b c) observer :=
  cocycle_condition a b c observer
\end{lstlisting}

\begin{interpretation}[A conservation law for meaning]
The cocycle is the ``conservation of energy'' for cultural meaning.
It says that the total emergence in any three-element composition is
\emph{independent of how you bracket the composition}.  Whether you
first compose $a$ with $b$ and then compose with $c$, or first compose
$b$ with $c$ and then compose with $a$, the \emph{total} emergence is
the same.

This is remarkable: the total meaning created by combining three
cultural elements is \emph{conserved} across different orderings of the
combination.  What changes is the \emph{distribution} of emergence
across steps, not the total.

The cocycle is what makes culture an \emph{algebra} rather than mere
chaos: it provides a structural constraint on how meaning is created,
ensuring that the creation of meaning, while nonlinear, is not arbitrary.
\end{interpretation}

\section{The Emergence-Weight Duality}

\begin{theorem}[Emergence-Weight Duality]
\label{thm:duality}
The two fundamental quantities---emergence and weight---are
related by:
\[
  \emerge(a, b, O)^2 \leq \rs(a \compose b, a \compose b) \cdot \rs(O, O)
  \qquad \text{and} \qquad
  \rs(a \compose b, a \compose b) \geq \rs(a, a).
\]
Emergence is bounded by weight, and weight is monotone under composition.
\end{theorem}

\begin{proof}
The first is axiom E3 (emergence boundedness).
The second is axiom E4 (compositional enrichment).
\end{proof}

\begin{lstlisting}
theorem emergence_weight_duality (a b obs : I) :
    (emergence a b obs) ^ 2 ≤
    rs (compose a b) (compose a b) * rs obs obs ∧
    rs (compose a b) (compose a b) ≥ rs a a :=
  ⟨emergence_bounded' a b obs, compose_enriches' a b⟩
\end{lstlisting}

\begin{remark}[Why this is surprising]
The duality says that emergence (creativity, surprise, nonlinearity)
and weight (substance, accumulation, persistence) are not independent.
They are bound together: more weight allows more emergence, and more
composition produces more weight.

This creates a \emph{positive feedback loop}: heavy cultural elements
can produce more emergence, which creates heavier elements, which can
produce even more emergence.  But the loop is bounded---emergence
\emph{squared} is bounded by weight, so the growth is at most
square-root-rate.  Culture can grow, but not exponentially.
\end{remark}

\section{The Linearity-Richness Dichotomy}

\begin{theorem}[Linearity Versus Richness]
\label{thm:dichotomy}
An element $a$ is left-linear if and only if its emergence is
universally zero:
\[
  \mathrm{leftLinear}(a) \iff \forall b, c.\; \emerge(a, b, c) = 0.
\]
Left-linear elements have zero semantic charge: $\sigma(a) = 0$.
\end{theorem}

\begin{proof}
Left-linearity is defined as universally zero emergence.
The semantic charge is $\emerge(a, a, a)$, which is zero by
left-linearity.
\end{proof}

\begin{lstlisting}
theorem linearity_richness_dichotomy (a : I) :
    leftLinear a → semanticCharge a = 0 := by
  intro h; unfold semanticCharge; exact h a a
\end{lstlisting}

\begin{interpretation}[The price of simplicity]
Left-linear elements are ``predictable'': their contribution to any
composition is purely additive.  They are the kitsch of the algebraic
world---no surprises, no emergence, no creative surplus.

Non-left-linear elements are ``rich'': they produce emergence, surprise,
creative interaction.  But they are also ``difficult''---their behavior
in compositions is nonlinear and hard to predict.

This creates a fundamental cultural dilemma: simplicity (left-linearity)
is manageable but sterile; richness (non-left-linearity) is creative
but unpredictable.  Every culture must navigate between these poles.
\end{interpretation}

\section{Culture Is Noncommutative}

\begin{theorem}[Noncommutativity of Culture]
\label{thm:noncommutative}
In any ideatic space with nontrivial emergence, there exist
elements where order matters:
\[
  \exist a, b, O.\;
  \rs(a \compose b, O) \neq \rs(b \compose a, O).
\]
Culture is generically noncommutative.
\end{theorem}

\begin{proof}
If $\emerge(a, b, O) \neq \emerge(b, a, O)$, then by the Order
Sensitivity Theorem (Chapter~\ref{ch:order}):
$\rs(a \compose b, O) - \rs(b \compose a, O) =
\emerge(a, b, O) - \emerge(b, a, O) \neq 0$.
\end{proof}

\begin{lstlisting}
theorem culture_noncommutative (a b observer : I)
    (h : emergence a b observer ≠ emergence b a observer) :
    rs (compose a b) observer ≠ rs (compose b a) observer := by
  intro heq
  apply h
  have h1 := rs_compose_eq a b observer
  have h2 := rs_compose_eq b a observer
  linarith
\end{lstlisting}

\begin{remark}[Why this is surprising]
Culture is \emph{not} a commutative algebra.  The order in which ideas
are composed matters.  ``Marx after Hegel'' is different from ``Hegel
after Marx.''  ``Cubism after Impressionism'' is different from
``Impressionism after Cubism.''

This is not merely a historical observation---it is a \emph{theorem}.
Any ideatic space with nontrivial emergence \emph{must} be
noncommutative.  Commutativity would require all emergence to be
symmetric ($\emerge(a,b,O) = \emerge(b,a,O)$ for all $a,b,O$), which
is a very strong condition that generically fails.

The noncommutativity of culture is the formal reason that \emph{history
matters}: the order of intellectual development is not interchangeable.
You cannot simply ``rearrange'' the history of ideas and get the same
result.
\end{remark}

\section{The Master Theorem: Cultural Thermodynamics}

We now state the culminating result of this volume: the
\textbf{Laws of Cultural Thermodynamics}.

\begin{theorem}[Cultural Thermodynamics]
\label{thm:cultural-thermo}
Any IdeaticSpace satisfies the following laws simultaneously:
\begin{enumerate}
\item \textbf{Conservation}: The cocycle
  $\emerge(a,b,O) + \emerge(a \compose b, c, O) =
   \emerge(b,c,O) + \emerge(a, b \compose c, O)$.
\item \textbf{Monotonicity}: Weight never decreases:
  $\rs(a \compose b, a \compose b) \geq \rs(a, a)$.
\item \textbf{Boundedness}: Emergence is bounded:
  $\emerge(a,b,O)^2 \leq \rs(a \compose b, a \compose b) \cdot \rs(O, O)$.
\item \textbf{Nondegeneracy}: Only void has zero weight:
  $\rs(a, a) = 0 \implies a = \void$.
\end{enumerate}
\end{theorem}

\begin{proof}
Laws (1)--(4) are respectively the cocycle condition (derived from
associativity), axiom E4, axiom E3, and axiom E2.  All are
consequences of the eight axioms of IdeaticSpace.
\end{proof}

\begin{lstlisting}
theorem cultural_thermodynamics (a b c observer : I) :
    (emergence a b observer +
     emergence (compose a b) c observer =
     emergence b c observer +
     emergence a (compose b c) observer) ∧
    (rs (compose a b) (compose a b) ≥ rs a a) ∧
    ((emergence a b observer) ^ 2 ≤
     rs (compose a b) (compose a b) *
     rs observer observer) ∧
    (rs a a = 0 → a = void) := by
  exact ⟨cocycle_condition a b c observer,
         compose_enriches' a b,
         emergence_bounded' a b observer,
         rs_self_zero_void a⟩
\end{lstlisting}

\begin{interpretation}[The four laws of cultural thermodynamics]
These four laws govern all cultural dynamics:

\textbf{First Law (Conservation)}: Total emergence is conserved across
different bracketings.  You can redistribute meaning creation across
steps, but you cannot create or destroy the total.  This is the
``conservation of creative energy.''

\textbf{Second Law (Monotonicity)}: Cultural weight never decreases.
Every composition adds substance.  Civilizations accumulate---they
cannot simplify without losing something.  This is the ``arrow of
cultural time.''

\textbf{Third Law (Boundedness)}: Creativity has a ceiling.  No
single act of composition can produce unlimited meaning.  Emergence
is bounded by the geometric mean of compositional weight and
observer weight.  This is the ``speed limit of innovation.''

\textbf{Fourth Law (Nondegeneracy)}: Only nothing is weightless.
Every genuine idea has positive self-resonance.  You cannot have
a non-trivial idea with zero substance.  This is the ``existence
of cultural mass.''

Together, these four laws describe a universe of meaning that is
rich (emergence is possible), constrained (emergence is bounded),
cumulative (weight grows), and grounded (only void is empty).
This is the algebra of culture.
\end{interpretation}

\bigskip

\begin{center}
\rule{0.5\textwidth}{0.5pt}
\end{center}

\bigskip

\noindent We have traveled from the decay of signs to the thermodynamics
of civilizations.  At every step, the same eight axioms---the axioms
of \emph{IdeaticSpace}---have produced theorems that formalize deep
intuitions from semiotics, aesthetics, narrative theory, translation
studies, and cultural criticism.

The theorems are often surprising, sometimes disturbing: meaning
necessarily decays (Chapter~\ref{ch:decay}), perfect translation is
impossible (Chapter~\ref{ch:translation}), novelty gets old
(Chapter~\ref{ch:defamiliarization}), beauty can destroy beauty
(Chapter~\ref{ch:kitsch}), parody creates originality
(Chapter~\ref{ch:parody}), and cultural memory grows heavier while
becoming less accessible (Chapter~\ref{ch:memory}).

But they are also, in their way, \emph{reassuring}: the laws of cultural
thermodynamics ensure that meaning creation is bounded but not
impossible, that cultural weight persists even when resonance fades,
and that the cocycle condition provides a structural guarantee that the
creation of meaning, while nonlinear, is not chaotic.

Culture, it turns out, is not a mystery.  It is an algebra.
And we have begun to learn its laws.


% ================================================================
\backmatter

\chapter*{Conclusion: The Laws of Cultural Thermodynamics}
\addcontentsline{toc}{chapter}{Conclusion}

We have arrived at a remarkable terminus.  From eight minimal axioms
about utterances, composition, and resonance, we have derived a
comprehensive theory of the \emph{paradoxes} of signs and culture.
The results organize into four ``laws of cultural thermodynamics'':

\begin{enumerate}
\item \textbf{The First Law} (Weight Conservation): Self-resonance never
  decreases under composition.  Culture accumulates; you cannot un-say
  something.  Every reading, every reinterpretation, every translation
  adds weight.

\item \textbf{The Second Law} (Emergence Boundedness): The emergent
  meaning $\emerge(a,b,c)$ cannot exceed
  $\sqrt{\rs(a \compose b, a \compose b) \cdot \rs(c,c)}$.  Meaning
  creation has an upper bound determined by the capacities of the parts.

\item \textbf{The Structural Law} (The Cocycle Condition): All cultural
  phenomena---sign chains, translation, defamiliarization, narrative,
  aesthetics---are constrained by the single identity
  \[
    \emerge(a,b,d) + \emerge(a \compose b, c, d)
    = \emerge(b,c,d) + \emerge(a, b \compose c, d).
  \]
  This is the ``conservation of emergence'' across levels of composition.

\item \textbf{The Third Law} (Nondegeneracy): Only silence ($\void$) has
  zero self-resonance.  Every genuine utterance has irreducible content.
\end{enumerate}

These four laws, all provable from the axioms, unify the 83 theorems of
this volume into a single coherent framework.  The paradoxes are not
anomalies---they are \emph{consequences} of the fundamental structure of
meaning.

Culture is non-commutative, emergence-bounded, cocycle-constrained, and
weight-monotonic.  These are not metaphors.  They are theorems.

\end{document}
